\documentclass[11pt,oneside]{book}

% ============================================================
% Packages
% ============================================================
\usepackage[a4paper,margin=1in]{geometry}
\usepackage[T1]{fontenc}
\usepackage[utf8]{inputenc}
\usepackage{lmodern}
\usepackage{microtype}
\usepackage{amsmath,amssymb,amsthm,mathtools}
\usepackage{bm}
\usepackage{enumitem}
\usepackage{booktabs,tabularx,array,multicol}
\usepackage{xcolor}
\usepackage{tikz}
\usetikzlibrary{arrows.meta,calc,patterns,positioning}
\usepackage{pgfplots}
\pgfplotsset{compat=1.18}
\usepackage{titlesec}
\usepackage{fancyhdr}
\usepackage{hyperref}
\usepackage{setspace}
\usepackage{needspace}

% ============================================================
% Colours and typography
% ============================================================
\definecolor{MainBlue}{HTML}{163A5F}
\definecolor{DeepBlue}{HTML}{0D2B45}
\definecolor{MainGreen}{HTML}{2F6B4F}
\definecolor{MainGold}{HTML}{9A6A16}
\definecolor{MainRed}{HTML}{8A3030}
\definecolor{Ink}{HTML}{222222}

\color{Ink}
\setstretch{1.08}
\setlength{\parindent}{0pt}
\setlength{\parskip}{4pt}
\setlist[itemize]{leftmargin=1.6em,itemsep=2pt,topsep=3pt}
\setlist[enumerate]{leftmargin=2em,itemsep=3pt,topsep=3pt}

% ============================================================
% Page style
% ============================================================
\pagestyle{fancy}
\fancyhf{}
\fancyhead[L]{\small\textsc{Undergraduate Calculus Notes}}
\fancyhead[R]{\small\thepage}
\fancyfoot{}
\renewcommand{\headrulewidth}{0.4pt}
\setlength{\headheight}{14pt}

\titleformat{\chapter}[display]
  {\normalfont\bfseries\color{DeepBlue}}
  {\filleft\Large\chaptername\ \thechapter}
  {1ex}
  {\titlerule\vspace{1.2ex}\Huge}
  [\vspace{1ex}\titlerule]

\titleformat{\section}
  {\Large\bfseries\color{MainBlue}}
  {\thesection}{0.7em}{}

\titleformat{\subsection}
  {\large\bfseries\color{MainGreen}}
  {\thesubsection}{0.7em}{}

% ============================================================
% Theorem environments
% ============================================================
\newtheoremstyle{studentthm}
{6pt}{6pt}{\itshape}{}
{\bfseries\color{DeepBlue}}{.}{0.5em}{}
\theoremstyle{studentthm}
\newtheorem{theorem}{Theorem}[chapter]
\newtheorem{proposition}[theorem]{Proposition}
\newtheorem{lemma}[theorem]{Lemma}
\newtheorem{corollary}[theorem]{Corollary}

\theoremstyle{definition}
\newtheorem{definition}[theorem]{Definition}
\newtheorem{example}[theorem]{Example}
\newtheorem{remark}[theorem]{Remark}

% ============================================================
% Commands
% ============================================================
\newcommand{\R}{\mathbb{R}}
\newcommand{\N}{\mathbb{N}}
\newcommand{\eps}{\varepsilon}
\newcommand{\abs}[1]{\left|#1\right|}
\newcommand{\norm}[1]{\left\lVert#1\right\rVert}
\newcommand{\pd}[2]{\frac{\partial #1}{\partial #2}}
\newcommand{\pdd}[3]{\frac{\partial^2 #1}{\partial #2\,\partial #3}}
\newcommand{\dd}{\,\mathrm d}

\hypersetup{
  colorlinks=true,
  linkcolor=MainBlue,
  urlcolor=MainGreen,
  pdfauthor={C. Zosangzuala},
  pdftitle={Functions of Two Variables: Limits, Partial Derivatives, Differentiability and Homogeneous Functions}
}

% ============================================================
\begin{document}

% ============================================================
% Cover
% ============================================================
\begin{titlepage}
\centering
\vspace*{2.0cm}

{\Huge\bfseries\color{DeepBlue}
Functions of Two Variables\par}

\vspace{0.35cm}

{\LARGE
Limits, Continuity, Partial Derivatives,\\
Differentiability, Composite Functions,\\
Change of Variables and Homogeneous Functions\par}

\vspace{1.25cm}

{\large
Extensive Lecture Notes for Undergraduate Calculus\par}

\vspace{0.9cm}

{\large
Real-valued functions of two variables; limits and repeated limits;
continuity; first- and second-order partial derivatives;
differentiability; Young's and Schwarz's theorems;
derivatives of composite functions; change of variables;
Jacobians; and Euler's theorem on homogeneous functions.\par}

\vfill

{\Large\bfseries C. Zosangzuala\par}
\vspace{0.25cm}
{\large Lecture Notes and Problem-Solving Companion\par}

\vspace{1cm}
\end{titlepage}

\frontmatter
\tableofcontents

% ============================================================
\chapter*{How to Use These Notes}
\addcontentsline{toc}{chapter}{How to Use These Notes}

Calculus of one variable studies a function such as
\[
y=f(x),
\]
where one real number determines the output. In this unit the input
consists of two real numbers:
\[
z=f(x,y).
\]
This apparently small change introduces several new ideas.

For a function of one variable, \(x\) approaches \(a\) essentially
from the left or from the right. In two variables, however,
\((x,y)\) can approach \((a,b)\) along infinitely many lines,
parabolas and other curves. Consequently, two-variable limits require
greater care.

Similarly, the derivative of a one-variable function gives one local
slope, whereas a surface \(z=f(x,y)\) has different rates of change
depending on how the point \((x,y)\) moves.

The conceptual order of the unit is
\[
\text{functions}
\longrightarrow
\text{limits and continuity}
\longrightarrow
\text{partial derivatives}
\longrightarrow
\text{differentiability}
\]
\[
\longrightarrow
\text{composite functions}
\longrightarrow
\text{change of variables}
\longrightarrow
\text{homogeneous functions}.
\]

The most important distinction to remember is
\[
\text{existence of partial derivatives}
\quad\not\Rightarrow\quad
\text{differentiability}.
\]

\mainmatter

% ============================================================
\chapter{Real-Valued Functions of Two Variables}
% ============================================================

\section*{Learning Objectives}

After this chapter a student should be able to:
\begin{itemize}
\item define a real-valued function of two variables;
\item find and sketch common domains in the \(xy\)-plane;
\item interpret the graph \(z=f(x,y)\) as a surface;
\item use level curves to understand a surface;
\item work with distance, neighbourhoods and deleted neighbourhoods
in \(\R^2\).
\end{itemize}

\section*{Historical Motivation}

The extension of calculus to several variables grew out of problems in
mechanics, astronomy, geometry and mathematical physics. Temperature,
pressure, potential energy and density naturally depend on more than
one independent variable. Euler, Lagrange, Clairaut, Jacobi and
Schwarz helped develop the language now used in multivariable
calculus.

\section{Definition}

\begin{definition}
Let \(D\subseteq\R^2\). A \emph{real-valued function of two real
variables} is a rule
\[
f:D\longrightarrow\R
\]
which assigns exactly one real number \(f(x,y)\) to each
\((x,y)\in D\).
\end{definition}

We usually write
\[
z=f(x,y).
\]

Here \(x,y\) are independent variables and \(z\) is the dependent
variable.

\paragraph{Everyday comparison.}
A quantity may depend on two independent inputs. For example, a taxi
fare may depend on both distance travelled and waiting time:
\[
\text{fare}=f(\text{distance},\text{waiting time}).
\]
Two inputs determine one output.

\begin{example}
For
\[
f(x,y)=x^2+y^2,
\]
every ordered pair is allowed, so
\[
\operatorname{Dom}(f)=\R^2.
\]
Since \(x^2+y^2\ge0\),
\[
\operatorname{Range}(f)=[0,\infty).
\]
\end{example}

\section{Finding Domains}

The familiar restrictions from one-variable calculus still apply:
\begin{itemize}
\item a denominator cannot be zero;
\item an even root requires a non-negative radicand;
\item a logarithm requires a positive argument.
\end{itemize}

\begin{example}
Find the domain of
\[
f(x,y)=\sqrt{9-x^2-y^2}.
\]
We require
\[
9-x^2-y^2\ge0,
\]
hence
\[
x^2+y^2\le9.
\]
The domain is the closed disc of radius \(3\) centred at the origin.
\end{example}

\begin{example}
Find the domain of
\[
f(x,y)=\log(4-x^2-y^2).
\]
We require
\[
4-x^2-y^2>0,
\]
hence
\[
x^2+y^2<4.
\]
The boundary circle is excluded.
\end{example}

\begin{example}
Find the domain of
\[
f(x,y)=\frac{x+y}{x-y}.
\]
The denominator must be nonzero:
\[
x-y\ne0.
\]
Thus the line \(y=x\) is removed from the plane.
\end{example}

\section{Graphs and Surfaces}

The graph of \(z=f(x,y)\) is
\[
\{(x,y,z)\in\R^3:z=f(x,y)\}.
\]
It is generally a surface.

For example,
\[
z=x^2+y^2
\]
is an upward-opening paraboloid.

\begin{center}
\begin{tikzpicture}
\begin{axis}[
width=10cm,height=7cm,
view={45}{27},
xlabel={$x$},ylabel={$y$},zlabel={$z$},
domain=-2:2,y domain=-2:2,
samples=23,samples y=23
]
\addplot3[surf,shader=interp,opacity=0.85]{x^2+y^2};
\end{axis}
\end{tikzpicture}
\end{center}

\section{Level Curves}

\begin{definition}
For a function \(z=f(x,y)\), a \emph{level curve} corresponding to
\(c\) is defined by
\[
f(x,y)=c.
\]
\end{definition}

For
\[
f(x,y)=x^2+y^2,
\]
the level curves are
\[
x^2+y^2=c.
\]
For \(c>0\), these are concentric circles.

\begin{center}
\begin{tikzpicture}[scale=0.9]
\draw[->] (-3.4,0)--(3.4,0) node[right] {$x$};
\draw[->] (0,-3.4)--(0,3.4) node[above] {$y$};
\draw[thick] (0,0) circle (1);
\draw[thick] (0,0) circle (1.8);
\draw[thick] (0,0) circle (2.6);
\node at (0.7,0.7) {$c=1$};
\node at (1.4,1.35) {$c=3.24$};
\node at (2.15,2.05) {$c=6.76$};
\end{tikzpicture}
\end{center}

Level curves are analogous to contour lines on a geographical map.


\section{From a Surface to Its Level Curves}

For the paraboloid
\[
z=x^2+y^2,
\]
a horizontal plane \(z=c\) cuts the surface in the curve
\[
x^2+y^2=c.
\]

Thus every horizontal slice of the surface produces a circle in the
\(xy\)-plane.

\begin{center}
\begin{tikzpicture}
\begin{axis}[
width=8.8cm,
height=6.2cm,
view={45}{26},
xlabel={$x$},
ylabel={$y$},
zlabel={$z$},
domain=-2:2,
y domain=-2:2,
samples=23,
samples y=23
]
\addplot3[surf,shader=interp,opacity=0.72]
{x^2+y^2};
\addplot3[very thick,domain=0:360,samples=80]
({sqrt(2)*cos(x)},{sqrt(2)*sin(x)},2);
\end{axis}
\end{tikzpicture}

\vspace{2mm}

\begin{tikzpicture}[scale=0.85]
\draw[->] (-3.2,0)--(3.2,0) node[right] {$x$};
\draw[->] (0,-3.2)--(0,3.2) node[above] {$y$};
\draw[thick] (0,0) circle (1);
\draw[thick] (0,0) circle (1.414);
\draw[thick] (0,0) circle (2);
\node at (1.0,0.35) {$c=1$};
\node at (1.55,0.75) {$c=2$};
\node at (2.25,1.1) {$c=4$};
\end{tikzpicture}
\end{center}

This is exactly analogous to contour lines on a topographical map:
each contour represents a fixed height.


\section{Distance and Neighbourhoods}

The distance from \(P=(x,y)\) to \(P_0=(a,b)\) is
\[
d(P,P_0)
=
\sqrt{(x-a)^2+(y-b)^2}.
\]

\begin{definition}
The \(\delta\)-neighbourhood of \((a,b)\) is
\[
N_\delta(a,b)
=
\left\{
(x,y):
\sqrt{(x-a)^2+(y-b)^2}<\delta
\right\}.
\]
\end{definition}

A deleted neighbourhood excludes the centre:
\[
0<
\sqrt{(x-a)^2+(y-b)^2}
<\delta.
\]

\paragraph{Neighbourhood analogy.}
Think of \((a,b)\) as a particular house. A neighbourhood contains
all points sufficiently close to the house. In the plane there are
infinitely many possible directions from which we may approach it.


\section{Interior, Boundary and Accumulation Points}

When limits and continuity are studied in two variables, it is useful to
know how a point sits relative to the domain.

\begin{definition}
A point \(P=(a,b)\) is an \emph{interior point} of a set \(D\subseteq\R^2\)
if there exists some \(\delta>0\) such that the whole disc
\[
N_\delta(a,b)
=
\{(x,y):\sqrt{(x-a)^2+(y-b)^2}<\delta\}
\]
lies inside \(D\).
\end{definition}

\begin{definition}
A point \(P=(a,b)\) is a \emph{boundary point} of \(D\) if every
neighbourhood of \(P\) contains points of \(D\) and also points which do
not belong to \(D\).
\end{definition}

\begin{definition}
A point \(P=(a,b)\) is an \emph{accumulation point} or \emph{limit point}
of \(D\) if every deleted neighbourhood
\[
0<\sqrt{(x-a)^2+(y-b)^2}<\delta
\]
contains at least one point of \(D\).
\end{definition}

These distinctions matter because the limit
\[
\lim_{(x,y)\to(a,b)}f(x,y)
\]
can be discussed whenever \((a,b)\) is an accumulation point of the
domain, even if \(f(a,b)\) itself is not defined.

\paragraph{Neighbourhood analogy.}
An interior house in a residential colony has a small circular region
around it that remains entirely inside the colony. A house on the
boundary does not have this property: every small neighbourhood around
it meets both the colony and the outside region. This is the geometric
idea behind interior and boundary points.

\section{Open, Closed and Bounded Regions}

A set \(D\subseteq\R^2\) is called \emph{open} if every point of \(D\) is
an interior point. For example,
\[
x^2+y^2<4
\]
is the open disc of radius \(2\).

A set is called \emph{closed} if it contains all its boundary points.
For example,
\[
x^2+y^2\le4
\]
is a closed disc.

A set is called \emph{bounded} if it is contained in some sufficiently
large disc:
\[
x^2+y^2<R^2
\]
for some \(R>0\).

Thus the closed disc
\[
x^2+y^2\le9
\]
is both closed and bounded, whereas the half-plane
\[
x+y>1
\]
is open but unbounded.

\section{Reading a Domain Geometrically}

Students should translate algebraic restrictions into plane regions.

\begin{center}
\begin{tabular}{ll}
Condition & Geometric meaning\\
\toprule
\(x^2+y^2<r^2\) & inside a circle, boundary excluded\\
\(x^2+y^2\le r^2\) & closed circular disc\\
\(y>x^2\) & above the parabola \(y=x^2\)\\
\(xy>0\) & first and third quadrants\\
\(xy<0\) & second and fourth quadrants\\
\(x\ne y\) & plane with the line \(y=x\) removed\\
\(x^2-y^2>0\) & \(|x|>|y|\)\\
\end{tabular}
\end{center}

\section{Additional Detailed Solved Problems}

\subsection*{Problem 4. A domain involving both a logarithm and a square root}

Find the domain of
\[
f(x,y)=\log\!\sqrt{9-x^2-y^2}.
\]

For the logarithm to be defined, its argument must be strictly positive.
Thus
\[
\sqrt{9-x^2-y^2}>0.
\]
This is equivalent to
\[
9-x^2-y^2>0,
\]
so
\[
x^2+y^2<9.
\]

Hence the domain is the open disc of radius \(3\). Notice that the
boundary circle is excluded, because on the boundary the square root is
zero and \(\log0\) is not defined.

\subsection*{Problem 5. A domain between two curves}

Find the domain of
\[
f(x,y)=\sqrt{y-x^2}\,\sqrt{4-y}.
\]

We require
\[
y-x^2\ge0
\]
and
\[
4-y\ge0.
\]
Thus
\[
x^2\le y\le4.
\]

The domain is the region above the parabola \(y=x^2\) and below the
horizontal line \(y=4\), including both boundary curves.

\subsection*{Problem 6. Level curves and interpretation}

Find and describe the level curves of
\[
f(x,y)=\frac{x^2}{4}+y^2.
\]

Putting \(f(x,y)=c\),
\[
\frac{x^2}{4}+y^2=c.
\]

If \(c<0\), there are no real level curves. If \(c=0\), the level set is
the single point \((0,0)\). If \(c>0\),
\[
\frac{x^2}{4c}+\frac{y^2}{c}=1.
\]
Hence the level curves are ellipses whose semiaxes are
\[
2\sqrt c
\quad\text{and}\quad
\sqrt c.
\]

As \(c\) increases, the ellipses expand outward. This gives a
two-dimensional picture of the three-dimensional paraboloid
\[
z=\frac{x^2}{4}+y^2.
\]



\section{How a Two-Variable Function Differs from a One-Variable Function}

For a function of one variable,
\[
y=f(x),
\]
a single input \(x\) determines the output. Its graph lies in the
two-dimensional \(xy\)-plane.

For a function of two variables,
\[
z=f(x,y),
\]
the input is an ordered pair \((x,y)\). Its graph therefore lies in
three-dimensional space.

This difference has several consequences.

\begin{enumerate}
\item A one-variable domain is usually an interval or a union of
intervals. A two-variable domain is usually a \emph{region} in the
plane.

\item A one-variable graph is a curve. A two-variable graph is usually
a surface.

\item In one variable, approaching \(a\) essentially means approaching
from the left or the right. In two variables, \((x,y)\) may approach
\((a,b)\) through infinitely many paths.

\item In one variable there is one ordinary derivative. In two
variables we may vary \(x\) while keeping \(y\) fixed, or vary \(y\)
while keeping \(x\) fixed. This produces partial derivatives.
\end{enumerate}

Thus the passage from one variable to two variables is not merely an
increase in notation; it changes the geometry of the theory.

\section{Examples from Geometry and Applications}

\subsection*{Height of a surface}

If the height of a hill above the \(xy\)-plane is given by
\[
z=100-2x^2-y^2,
\]
then each ground position \((x,y)\) determines exactly one height
\(z\).

\subsection*{Temperature on a plate}

If
\[
T(x,y)=25+3x-2y,
\]
then \(T(x,y)\) represents the temperature at the point \((x,y)\) on a
thin plate.

\subsection*{Distance from the origin}

The function
\[
f(x,y)=\sqrt{x^2+y^2}
\]
assigns to every point \((x,y)\) its Euclidean distance from the
origin.

\subsection*{Potential-type functions}

Expressions such as
\[
V(x,y)=\frac{1}{\sqrt{x^2+y^2}}
\]
occur naturally in physical models. Its domain excludes the origin,
because the denominator vanishes there.

These examples show why domains should always be determined before
performing differentiation or limit calculations.

\section{A Systematic Strategy for Finding Domains}

When the formula contains several restrictions, it is helpful to
analyze them one at a time.

\begin{enumerate}
\item Write the restriction imposed by every denominator.

\item Write the restriction imposed by every square root or other even
root.

\item Write the restriction imposed by every logarithm.

\item Intersect all the resulting regions.

\item Translate the final inequalities into a geometric description.
\end{enumerate}

\begin{example}
Find the domain of
\[
f(x,y)
=
\frac{\log(y-x^2)}{\sqrt{9-x^2-y^2}}.
\]

For the logarithm,
\[
y-x^2>0,
\]
so
\[
y>x^2.
\]

For the denominator involving the square root, we need
\[
9-x^2-y^2>0,
\]
not merely \(\ge0\), because the square root appears in the
denominator.

Thus
\[
x^2+y^2<9.
\]

Therefore the domain is the part of the open disc
\[
x^2+y^2<9
\]
which lies strictly above the parabola
\[
y=x^2.
\]
\end{example}

\section{Vertical Traces and Horizontal Traces}

The graph of a surface can often be understood by taking cross-sections.

For
\[
z=f(x,y),
\]
fixing \(y=b\) gives the vertical trace
\[
z=f(x,b).
\]

Fixing \(x=a\) gives
\[
z=f(a,y).
\]

Fixing \(z=c\) gives the horizontal trace or level curve
\[
f(x,y)=c.
\]

For example, let
\[
z=x^2+y^2.
\]

If \(y=0\),
\[
z=x^2,
\]
a parabola in the \(xz\)-plane.

If \(x=0\),
\[
z=y^2,
\]
a parabola in the \(yz\)-plane.

If \(z=c>0\),
\[
x^2+y^2=c,
\]
a circle.

These three families of traces together make the shape of the
paraboloid much easier to visualize.

\section{More Detailed Problems on Domains and Level Curves}

\subsection*{Problem 7. Domain involving a quotient and logarithm}

Find the domain of
\[
f(x,y)
=
\log\left(\frac{x+y}{x-y}\right).
\]

The logarithm requires
\[
\frac{x+y}{x-y}>0.
\]

A quotient is positive when numerator and denominator have the same
sign. Hence either
\[
x+y>0,\qquad x-y>0,
\]
or
\[
x+y<0,\qquad x-y<0.
\]

These two cases describe two opposite angular regions bounded by the
lines
\[
y=x,\qquad y=-x.
\]

The boundary lines are excluded.

\subsection*{Problem 8. A family of level curves}

Let
\[
f(x,y)=xy.
\]

The level curves satisfy
\[
xy=c.
\]

If \(c\ne0\),
\[
y=\frac{c}{x},
\]
which gives a family of rectangular hyperbolas.

If \(c=0\),
\[
xy=0,
\]
so the level set consists of the two coordinate axes.

Thus a single surface can produce qualitatively different level sets
for different values of \(c\).

\subsection*{Problem 9. Comparing two surfaces}

Compare
\[
z=x^2+y^2
\]
with
\[
z=x^2-y^2.
\]

For \(z=x^2+y^2\), both coordinate traces are upward-opening
parabolas. Its level curves are circles, so the surface is an elliptic
paraboloid.

For \(z=x^2-y^2\), the trace \(y=0\) is
\[
z=x^2,
\]
while the trace \(x=0\) is
\[
z=-y^2.
\]

Thus one direction curves upward and the perpendicular direction
curves downward. The surface is a saddle. Its level curves
\[
x^2-y^2=c
\]
are hyperbolas for \(c\ne0\).


\section{Detailed Solved Problems}

\subsection*{Problem 1. Domain of a square-root function}
Find the domain of
\[
f(x,y)=\sqrt{16-x^2-4y^2}.
\]

We require
\[
16-x^2-4y^2\ge0.
\]
Thus
\[
x^2+4y^2\le16.
\]
Dividing by \(16\),
\[
\frac{x^2}{16}+\frac{y^2}{4}\le1.
\]
The domain is the closed elliptical region inside and on the ellipse.

\subsection*{Problem 2. Domain of a logarithmic function}
Find the domain of
\[
f(x,y)=\log(x-y^2).
\]
We require
\[
x-y^2>0,
\]
so
\[
x>y^2.
\]
The domain is the region to the right of the parabola \(x=y^2\);
the parabola itself is excluded.

\subsection*{Problem 3. Level curves}
Describe the level curves of
\[
f(x,y)=x^2+4y^2.
\]
Setting \(f(x,y)=c\),
\[
x^2+4y^2=c.
\]
For \(c>0\), these are ellipses centred at the origin.

\section{Graded Practice Problems}

\subsection*{Level I: Foundations}
\begin{enumerate}
\item Find the domain of \(x^2-3xy+y^2\).
\item Find the domain of \(\sqrt{25-x^2-y^2}\).
\item Find the domain of \(1/(x+y)\).
\item Describe the level curves of \(f(x,y)=x+y\).
\end{enumerate}

\subsection*{Level II: Standard Problems}
\begin{enumerate}
\item Find and sketch the domain of \(\sqrt{y-x^2}\).
\item Find the domain of \(\log(9-x^2-y^2)\).
\item Describe the level curves of \(4x^2+y^2\).
\item Explain geometrically the surface \(z=4-x^2-y^2\).
\end{enumerate}

\subsection*{Level III: Challenging}
\begin{enumerate}
\item Determine the domain of
\[
\sqrt{\frac{x-y}{x+y}}.
\]
\item Describe the level curves of \(x/y\).
\item Compare the surfaces \(z=x^2+y^2\) and \(z=x^2-y^2\).
\end{enumerate}

% ============================================================
\chapter{Limits, Repeated Limits and Continuity}
% ============================================================

\section*{Learning Objectives}
Students should be able to:
\begin{itemize}
\item state and use the \(\eps\)-\(\delta\) definition of a
two-variable limit;
\item test nonexistence by different paths;
\item understand why a few agreeing paths do not prove existence;
\item use inequalities and polar coordinates to prove existence;
\item compute repeated limits;
\item distinguish repeated limits from a joint limit;
\item test continuity.
\end{itemize}

\section*{Historical Motivation}

The rigorous ideas of Cauchy and Weierstrass extend naturally to the
plane, but the geometry changes. A point in the plane can be
approached in infinitely many ways. A correct definition must
therefore control all approaches simultaneously.

\section{Definition of a Joint Limit}

\begin{definition}
We say
\[
\lim_{(x,y)\to(a,b)}f(x,y)=L
\]
if, for every \(\eps>0\), there exists \(\delta>0\) such that
\[
0<
\sqrt{(x-a)^2+(y-b)^2}
<\delta
\]
implies
\[
\abs{f(x,y)-L}<\eps.
\]
\end{definition}

The distance condition does not prescribe a path, so the definition
controls every possible approach to \((a,b)\).

\section{Visual Meaning of \(\eps\) and \(\delta\)}

\begin{center}
\begin{tikzpicture}[scale=1.0]
\draw[->] (-0.5,0)--(6,0) node[right] {$x$};
\draw[->] (0,-0.5)--(0,5) node[above] {$y$};
\coordinate (A) at (3,2.2);
\draw[thick] (A) circle (1.35);
\fill (A) circle (2pt);
\node[below right] at (A) {$(a,b)$};
\draw[<->] (3,2.2)--(4.35,2.2);
\node[above] at (3.7,2.2) {$\delta$};
\node at (3,4.1) {$N_\delta(a,b)$};
\end{tikzpicture}
\end{center}

The \(\delta\)-disc controls the inputs. The \(\eps\)-condition controls
the outputs.

\section{A Direct \(\eps\)-\(\delta\) Proof}

\begin{example}
Prove
\[
\lim_{(x,y)\to(0,0)}(3x-2y)=0.
\]
\end{example}

\textbf{Proof.}
We want to make
\[
\abs{3x-2y}
\]
smaller than a prescribed \(\eps>0\).

By the triangle inequality,
\[
\abs{3x-2y}
\le3\abs{x}+2\abs{y}.
\]

Also,
\[
\abs{x}\le\sqrt{x^2+y^2},
\qquad
\abs{y}\le\sqrt{x^2+y^2}.
\]
Hence
\[
\abs{3x-2y}
\le5\sqrt{x^2+y^2}.
\]

Choose
\[
\delta=\frac{\eps}{5}.
\]
Then
\[
0<\sqrt{x^2+y^2}<\delta
\]
implies
\[
\abs{3x-2y}
<
5\delta
=
\eps.
\]
Therefore
\[
\lim_{(x,y)\to(0,0)}(3x-2y)=0.
\]

\section{Standard Limit Laws}

Suppose
\[
\lim_{(x,y)\to(a,b)}f(x,y)=L
\]
and
\[
\lim_{(x,y)\to(a,b)}g(x,y)=M.
\]
Then:
\[
\lim(f+g)=L+M,
\]
\[
\lim(f-g)=L-M,
\]
\[
\lim(fg)=LM,
\]
and, if \(M\ne0\),
\[
\lim\frac{f}{g}=\frac{L}{M}.
\]

The proofs follow from the same \(\eps\)-\(\delta\) ideas used in
one-variable calculus.


\section{Visualizing Different Paths of Approach}

The same point can be approached along infinitely many curves. The
figure below shows several possible approaches to the origin.

\begin{center}
\begin{tikzpicture}[scale=1.05,>=Stealth]
\draw[->] (-3.3,0)--(3.3,0) node[right] {$x$};
\draw[->] (0,-2.7)--(0,2.9) node[above] {$y$};

\fill (0,0) circle (2pt) node[below right] {$(0,0)$};

\draw[thick,->] (-3,-1.5)--(-0.18,-0.09);
\node[below] at (-2.1,-1.05) {$y=\frac12x$};

\draw[thick,->] (2.8,-1.4)--(0.18,-0.09);
\node[below] at (2.05,-1.05) {$y=-\frac12x$};

\draw[thick,domain=-2.2:-0.16,samples=40,->]
plot (\x,{0.35*(\x)^2});
\node[left] at (-1.75,1.15) {$y=0.35x^2$};

\draw[thick,domain=2.2:0.16,samples=40,->]
plot (\x,{0.12*(\x)^3});
\node[right] at (1.6,0.7) {$y=0.12x^3$};

\draw[dashed] (0,0) circle (1.8);
\node at (-2.45,2.1) {many possible paths};
\end{tikzpicture}
\end{center}

A genuine joint limit must give the same limiting value along
\emph{all} of these approaches and along every other admissible path.


\section{Path Testing}

If the joint limit exists, every path to the point must give the same
value. Thus two paths with different limiting values prove
nonexistence immediately.

\begin{example}
Examine
\[
\lim_{(x,y)\to(0,0)}
\frac{x^2-y^2}{x^2+y^2}.
\]

Along \(y=0\),
\[
\frac{x^2}{x^2}=1.
\]
Along \(x=0\),
\[
\frac{-y^2}{y^2}=-1.
\]
The two limits are different. Therefore the joint limit does not
exist.
\end{example}

\section{Why Agreement Along Straight Lines Is Not Enough}

Consider
\[
f(x,y)
=
\frac{x^2y}{x^4+y^2}.
\]
Along the line \(y=mx\),
\[
f(x,mx)
=
\frac{mx^3}{x^4+m^2x^2}
=
\frac{mx}{x^2+m^2}
\longrightarrow0.
\]
Thus every fixed straight line gives \(0\).

But along the parabola \(y=x^2\),
\[
f(x,x^2)
=
\frac{x^4}{x^4+x^4}
=
\frac12.
\]
Hence the joint limit does not exist.

\paragraph{Key principle.}
Different paths are excellent for proving nonexistence. Agreement
along selected paths is not, by itself, a proof of existence.

\section{Proving Existence by an Inequality}

Consider
\[
f(x,y)=\frac{x^2y}{x^2+y^2}.
\]
Then
\[
\abs{f(x,y)}
=
\abs{y}\frac{x^2}{x^2+y^2}
\le\abs{y}.
\]
Since
\[
\abs{y}\to0
\quad\text{as}\quad
(x,y)\to(0,0),
\]
the Squeeze (Sandwich) Theorem gives
\[
\lim_{(x,y)\to(0,0)}
\frac{x^2y}{x^2+y^2}
=0.
\]

\section{Polar Coordinates in Limit Problems}

At the origin, set
\[
x=r\cos\theta,\qquad y=r\sin\theta.
\]
Then
\[
x^2+y^2=r^2
\]
and
\[
(x,y)\to(0,0)
\quad\Longleftrightarrow\quad
r\to0.
\]

For example,
\[
\frac{x^2y}{x^2+y^2}
=
r\cos^2\theta\sin\theta.
\]
Hence
\[
\abs{r\cos^2\theta\sin\theta}\le r,
\]
so the limit is \(0\).

\section{Repeated Limits}

\begin{definition}
The repeated limit
\[
\lim_{x\to a}
\left[
\lim_{y\to b}f(x,y)
\right]
\]
is obtained by first holding \(x\) fixed and taking \(y\to b\), then
taking \(x\to a\).
\end{definition}

The reverse order is
\[
\lim_{y\to b}
\left[
\lim_{x\to a}f(x,y)
\right].
\]

\begin{example}
For
\[
f(x,y)=\frac{x^2-y^2}{x^2+y^2},
\]
fix \(x\ne0\). Then
\[
\lim_{y\to0}f(x,y)=1.
\]
Thus
\[
\lim_{x\to0}\lim_{y\to0}f(x,y)=1.
\]

On the other hand, fixing \(y\ne0\),
\[
\lim_{x\to0}f(x,y)=-1,
\]
so
\[
\lim_{y\to0}\lim_{x\to0}f(x,y)=-1.
\]
The repeated limits are different.
\end{example}

\section{Relation Between Joint and Repeated Limits}

\begin{theorem}
If the joint limit
\[
\lim_{(x,y)\to(a,b)}f(x,y)=L
\]
exists and the relevant inner limits exist in a neighbourhood, then
both repeated limits, when defined, must equal \(L\).
\end{theorem}

However, the converse is false.

\begin{example}
Let
\[
f(x,y)=
\begin{cases}
\dfrac{xy}{x^2+y^2},&(x,y)\ne(0,0),\\[1ex]
0,&(x,y)=(0,0).
\end{cases}
\]
For fixed \(x\),
\[
\lim_{y\to0}f(x,y)=0,
\]
so
\[
\lim_{x\to0}\lim_{y\to0}f(x,y)=0.
\]
Similarly,
\[
\lim_{y\to0}\lim_{x\to0}f(x,y)=0.
\]

But along \(y=x\),
\[
f(x,x)=\frac12.
\]
Therefore the joint limit does not exist.

Thus
\[
\text{equal repeated limits}
\not\Rightarrow
\text{existence of the joint limit}.
\]
\end{example}

\section{Continuity}

\begin{definition}
A function \(f(x,y)\) is continuous at \((a,b)\) if:
\begin{enumerate}
\item \(f(a,b)\) is defined;
\item \(\lim_{(x,y)\to(a,b)}f(x,y)\) exists;
\item
\[
\lim_{(x,y)\to(a,b)}f(x,y)=f(a,b).
\]
\end{enumerate}
\end{definition}

Polynomials in \(x,y\) are continuous everywhere. Rational functions
are continuous wherever their denominators are nonzero.

\begin{example}
Define
\[
f(x,y)=
\begin{cases}
\dfrac{x^2y}{x^2+y^2},&(x,y)\ne(0,0),\\[1ex]
0,&(x,y)=(0,0).
\end{cases}
\]
Since
\[
\abs{\frac{x^2y}{x^2+y^2}}
\le\abs{y}\to0,
\]
we have
\[
\lim_{(x,y)\to(0,0)}f(x,y)=0=f(0,0).
\]
Hence \(f\) is continuous at the origin.
\end{example}


\section{A Necessary Path Condition for a Joint Limit}

\begin{theorem}
Suppose
\[
\lim_{(x,y)\to(a,b)}f(x,y)=L.
\]
Let
\[
x=\phi(t),\qquad y=\psi(t)
\]
be any path for which
\[
(\phi(t),\psi(t))\to(a,b)
\]
as \(t\to t_0\), with \((\phi(t),\psi(t))\ne(a,b)\) near \(t_0\).
Then
\[
\lim_{t\to t_0}f(\phi(t),\psi(t))=L.
\]
\end{theorem}

\begin{proof}
Because the joint limit equals \(L\), for every \(\eps>0\) there exists
\(\delta>0\) such that
\[
0<
\sqrt{(x-a)^2+(y-b)^2}
<\delta
\]
implies
\[
|f(x,y)-L|<\eps.
\]

Since
\[
(\phi(t),\psi(t))\to(a,b),
\]
there exists a neighbourhood of \(t_0\) in which
\[
0<
\sqrt{
(\phi(t)-a)^2+(\psi(t)-b)^2
}
<\delta.
\]

Therefore
\[
|f(\phi(t),\psi(t))-L|<\eps.
\]
Hence the limit along the path is also \(L\).
\end{proof}

The theorem gives an extremely useful logical principle:

\[
\text{joint limit exists}
\quad\Longrightarrow\quad
\text{every admissible path gives the same limit}.
\]

Therefore, to prove nonexistence, it is enough to find two paths that
give different answers.

\section{The Sequential Criterion for a Two-Variable Limit}

There is another useful way to understand a joint limit.

\begin{theorem}[Sequential criterion]
Let \((a,b)\) be an accumulation point of the domain of \(f\). Then
\[
\lim_{(x,y)\to(a,b)}f(x,y)=L
\]
if and only if for every sequence of points
\[
(x_n,y_n)\ne(a,b)
\]
satisfying
\[
(x_n,y_n)\to(a,b),
\]
we have
\[
f(x_n,y_n)\to L.
\]
\end{theorem}

This theorem says that a two-variable limit must survive \emph{every
possible sequence of points} approaching the point.

For nonexistence, it is enough to find two sequences approaching the
same point but producing different limiting values.

\begin{example}
Let
\[
f(x,y)=\frac{x^2-y^2}{x^2+y^2}.
\]
Choose
\[
(x_n,y_n)=\left(\frac1n,0\right).
\]
Then
\[
f(x_n,y_n)=1.
\]
But for
\[
(u_n,v_n)=\left(0,\frac1n\right),
\]
\[
f(u_n,v_n)=-1.
\]
Both sequences approach \((0,0)\), but the function values approach
different limits. Hence the joint limit does not exist.
\end{example}

\section{How to Choose a Useful Path}

There is no mechanical rule that always finds the right path, but the
algebra often suggests one.

\begin{enumerate}
\item If the denominator contains \(x^2+y^2\), try
\[
y=mx.
\]

\item If the denominator contains \(x^4+y^2\), try to balance the
powers by choosing
\[
y=kx^2.
\]

\item If the denominator contains \(x^6+y^2\), a natural curved path is
\[
y=kx^3.
\]

\item More generally, if two terms have different powers, choose a path
which makes them of comparable order.
\end{enumerate}

\begin{example}
Consider
\[
f(x,y)=\frac{x^2y}{x^4+y^2}.
\]
The denominator contains \(x^4\) and \(y^2\). To make these of the same
order, choose
\[
y=kx^2.
\]
Then
\[
f(x,kx^2)
=
\frac{kx^4}{x^4+k^2x^4}
=
\frac{k}{1+k^2}.
\]
The answer depends on \(k\). Therefore the limit cannot exist.
\end{example}

\section{Homogeneous Expressions and Polar Coordinates}

Suppose the numerator and denominator are homogeneous expressions.
Polar coordinates frequently expose the exact power of \(r\) which
controls the limit.

For example,
\[
f(x,y)
=
\frac{x^3y^2}{(x^2+y^2)^2}.
\]
Put
\[
x=r\cos\theta,\qquad y=r\sin\theta.
\]
Then
\[
x^3y^2
=
r^5\cos^3\theta\sin^2\theta
\]
and
\[
(x^2+y^2)^2=r^4.
\]
Hence
\[
f(x,y)
=
r\cos^3\theta\sin^2\theta.
\]
Since
\[
|\cos^3\theta\sin^2\theta|\le1,
\]
we obtain
\[
|f(x,y)|\le r.
\]
Therefore
\[
\lim_{(x,y)\to(0,0)}f(x,y)=0.
\]

The useful pattern is:

\[
r^\alpha\times\text{bounded angular factor}.
\]

If \(\alpha>0\), the limit is usually \(0\). If \(\alpha=0\), the
angular factor may reveal path dependence. If \(\alpha<0\), the
expression may become unbounded.

\section{Repeated Limits: What They Do and Do Not Tell Us}

Repeated limits inspect the variables in a prescribed order:
\[
\lim_{x\to a}\lim_{y\to b}f(x,y),
\qquad
\lim_{y\to b}\lim_{x\to a}f(x,y).
\]

The logical relations should be remembered carefully.

If the joint limit exists under the usual local hypotheses, then any
repeated limits which are defined must agree with that joint limit.

But the converse fails:

\[
\text{both repeated limits exist and are equal}
\quad\not\Rightarrow\quad
\text{joint limit exists}.
\]

\begin{example}
For
\[
f(x,y)=\frac{xy}{x^2+y^2},
\]
both repeated limits at the origin are zero, because fixing one
variable and sending the other to zero produces zero. However, along
\(y=x\),
\[
f(x,x)=\frac12.
\]
Hence the joint limit does not exist.
\end{example}

\section{Continuity: Algebraic Theorems}

\begin{theorem}
If \(f\) and \(g\) are continuous at \((a,b)\), then
\[
f+g,\qquad f-g,\qquad fg
\]
are continuous at \((a,b)\). If
\[
g(a,b)\ne0,
\]
then
\[
\frac{f}{g}
\]
is also continuous at \((a,b)\).
\end{theorem}

\begin{theorem}[Continuity of a composite]
Suppose \(u=g(x,y)\) is continuous at \((a,b)\), and \(\phi\) is a
one-variable function continuous at \(g(a,b)\). Then
\[
\phi(g(x,y))
\]
is continuous at \((a,b)\).
\end{theorem}

Thus functions such as
\[
e^{x^2+y^2},
\qquad
\sin(xy),
\qquad
\sqrt{1+x^2+y^2}
\]
are continuous everywhere.

\section{Additional Detailed Solved Problems}

\subsection*{Problem 5. Prove a limit by a direct estimate}

Prove
\[
\lim_{(x,y)\to(0,0)}
\frac{x^2y^2}{x^2+y^2}
=0.
\]

We use
\[
x^2\le x^2+y^2,
\qquad
y^2\le x^2+y^2.
\]
Therefore
\[
x^2y^2
\le
(x^2+y^2)^2.
\]
Hence
\[
0\le
\frac{x^2y^2}{x^2+y^2}
\le
x^2+y^2.
\]
As
\[
(x,y)\to(0,0),
\]
the right-hand side tends to zero. Therefore, by the Sandwich Theorem,
\[
\lim_{(x,y)\to(0,0)}
\frac{x^2y^2}{x^2+y^2}
=0.
\]

\subsection*{Problem 6. A curved-path counterexample}

Examine
\[
\lim_{(x,y)\to(0,0)}
\frac{x^2y}{x^4+y^2}.
\]

Along \(y=0\),
\[
f(x,0)=0.
\]

Along the curved path
\[
y=x^2,
\]
\[
f(x,x^2)
=
\frac{x^4}{x^4+x^4}
=
\frac12.
\]

The two path limits are different. Hence the joint limit does not
exist.

\subsection*{Problem 7. Repeated limits can agree while the joint limit fails}

Let
\[
f(x,y)
=
\frac{xy}{x^2+y^2}.
\]

For fixed \(x\ne0\),
\[
\lim_{y\to0}f(x,y)=0,
\]
so
\[
\lim_{x\to0}\lim_{y\to0}f(x,y)=0.
\]

Similarly,
\[
\lim_{y\to0}\lim_{x\to0}f(x,y)=0.
\]

However, along \(y=x\),
\[
f(x,x)=\frac12.
\]

Therefore the joint limit does not exist even though the repeated
limits exist and are equal.

\subsection*{Problem 8. Continuity at a specially defined point}

Define
\[
f(x,y)
=
\begin{cases}
\dfrac{x^3+y^3}{x^2+y^2},&(x,y)\ne(0,0),\\[1ex]
0,&(x,y)=(0,0).
\end{cases}
\]

We estimate
\[
|x^3+y^3|
\le
|x|^3+|y|^3.
\]
Let
\[
r=\sqrt{x^2+y^2}.
\]
Then
\[
|x|\le r,\qquad |y|\le r,
\]
so
\[
|x^3+y^3|\le2r^3.
\]
Since
\[
x^2+y^2=r^2,
\]
\[
|f(x,y)|
\le2r.
\]
As \(r\to0\),
\[
f(x,y)\to0=f(0,0).
\]
Hence \(f\) is continuous at the origin.



\section{Which Method Should I Use for a Two-Variable Limit?}

A useful first decision is whether we are trying to \emph{prove
existence} or \emph{prove nonexistence}. The following guide helps
students choose an efficient method.

\begin{center}
\renewcommand{\arraystretch}{1.35}
\begin{tabularx}{0.94\textwidth}{>{\bfseries}p{0.31\textwidth}X}
\toprule
Situation & Recommended first method\\
\midrule
Direct substitution gives an ordinary real number
&
Use continuity and substitute directly.
\\

You suspect that the limit does not exist
&
Compare two convenient paths. Start with \(x=0\), \(y=0\), or
\(y=mx\), then try a curved path if necessary.
\\

The denominator contains \(x^2+y^2\)
&
Try an estimate involving
\(\sqrt{x^2+y^2}\), or use polar coordinates.
\\

Different powers such as \(x^4+y^2\) occur
&
Choose a balancing path such as \(y=kx^2\).
\\

Several tested paths give the same value
&
Do not conclude that the limit exists. Seek a uniform estimate,
polar-coordinate argument, or an \(\eps\)-\(\delta\) proof.
\\

A rigorous proof of existence is required
&
Use an inequality, the Sandwich Theorem, polar coordinates with a
uniform angular bound, or the \(\eps\)-\(\delta\) definition.
\\

Repeated limits are requested
&
Take one variable to its limit while holding the other fixed, then
take the remaining limit. Repeat in the opposite order.
\\
\bottomrule
\end{tabularx}
\end{center}

The most important distinction is:

\[
\text{different path values}
\Longrightarrow
\text{joint limit does not exist},
\]
whereas
\[
\text{same values on several paths}
\]
does not by itself prove existence.



\section{Why the Two-Variable Limit Definition Uses Distance}

In one variable,
\[
|x-a|<\delta
\]
means that \(x\) lies inside a small interval centered at \(a\).

In two variables, the natural analogue is
\[
\sqrt{(x-a)^2+(y-b)^2}<\delta.
\]

This describes a disc centered at \((a,b)\). The condition is
direction-free: it does not matter whether the point approaches from
the left, right, above, below, diagonally, or along a curved path.

That is exactly what a joint limit must control.

\section{A Second Direct \(\eps\)-\(\delta\) Proof}

Prove that
\[
\lim_{(x,y)\to(0,0)}(x^2+y^2)=0.
\]

Let \(\eps>0\) be given. We want
\[
|x^2+y^2|<\eps.
\]

But
\[
x^2+y^2
=
\left(\sqrt{x^2+y^2}\right)^2.
\]

Thus if
\[
\sqrt{x^2+y^2}<\delta,
\]
then
\[
x^2+y^2<\delta^2.
\]

Choose
\[
\delta=\sqrt{\eps}.
\]

Then
\[
0<\sqrt{x^2+y^2}<\delta
\]
implies
\[
|x^2+y^2|
=
x^2+y^2
<
\delta^2
=
\eps.
\]

Therefore
\[
\lim_{(x,y)\to(0,0)}(x^2+y^2)=0.
\]

\section{A More Involved \(\eps\)-\(\delta\) Proof}

Prove
\[
\lim_{(x,y)\to(1,2)}(x+2y)=5.
\]

We have
\[
\begin{aligned}
|(x+2y)-5|
&=
|(x-1)+2(y-2)|\\
&\le
|x-1|+2|y-2|.
\end{aligned}
\]

Also,
\[
|x-1|
\le
\sqrt{(x-1)^2+(y-2)^2},
\]
and
\[
|y-2|
\le
\sqrt{(x-1)^2+(y-2)^2}.
\]

Therefore
\[
|(x+2y)-5|
\le
3\sqrt{(x-1)^2+(y-2)^2}.
\]

Choose
\[
\delta=\frac{\eps}{3}.
\]

Then
\[
0<
\sqrt{(x-1)^2+(y-2)^2}
<\delta
\]
implies
\[
|(x+2y)-5|<3\delta=\eps.
\]

Hence
\[
\lim_{(x,y)\to(1,2)}(x+2y)=5.
\]

\section{Limit Laws}

The familiar algebra of one-variable limits remains valid.

Suppose
\[
\lim_{(x,y)\to(a,b)}f(x,y)=L
\]
and
\[
\lim_{(x,y)\to(a,b)}g(x,y)=M.
\]

Then
\[
\lim(f+g)=L+M,
\]
\[
\lim(f-g)=L-M,
\]
\[
\lim(fg)=LM,
\]
and if \(M\ne0\),
\[
\lim\frac{f}{g}=\frac{L}{M}.
\]

Also, for a constant \(c\),
\[
\lim cf=cL.
\]

These laws allow direct substitution into polynomials and rational
functions whenever the denominator does not vanish.

\subsection*{Why direct substitution works for polynomials}

A polynomial in \(x,y\) is built from constants, coordinate functions,
sums and products. Since
\[
\lim_{(x,y)\to(a,b)}x=a,
\qquad
\lim_{(x,y)\to(a,b)}y=b,
\]
the limit laws imply
\[
\lim_{(x,y)\to(a,b)}P(x,y)=P(a,b)
\]
for every polynomial \(P\).

\section{Path Tests: A Hierarchy of Increasing Strength}

When nonexistence is suspected, try paths in stages.

\subsection*{Stage 1: Coordinate axes}

Try
\[
x=0
\quad\text{and}\quad
y=0.
\]

If these give different values, stop: the limit does not exist.

\subsection*{Stage 2: Straight lines}

Try
\[
y=mx.
\]

If the resulting limit depends on \(m\), the joint limit does not
exist.

\subsection*{Stage 3: Curved paths}

If all straight lines give the same value, look at the powers in the
formula and choose a balancing path such as
\[
y=kx^2,\qquad y=kx^3,\qquad x=ky^2.
\]

\subsection*{Stage 4: Uniform proof}

If no path reveals a contradiction and the limit appears to exist,
path checking should stop. Replace it by a uniform estimate or polar
argument.

This is an important change in strategy:
\[
\text{paths are best for disproving,}
\]
while
\[
\text{inequalities are usually better for proving existence.}
\]

\section{A General Polar-Coordinate Test}

Suppose that after substituting
\[
x=r\cos\theta,\qquad y=r\sin\theta,
\]
a function becomes
\[
f(x,y)=r^\alpha \Phi(\theta),
\]
where
\[
|\Phi(\theta)|\le M
\]
for all \(\theta\).

If \(\alpha>0\), then
\[
|f(x,y)|
\le
Mr^\alpha
\to0.
\]

Hence
\[
\lim_{(x,y)\to(0,0)}f(x,y)=0.
\]

The boundedness of the angular factor is essential. Merely rewriting in
polar coordinates is not itself a proof.

\section{A Polar Example with a Nonzero Angular Dependence}

Consider
\[
f(x,y)
=
\frac{x^2-y^2}{x^2+y^2}.
\]

With
\[
x=r\cos\theta,\qquad y=r\sin\theta,
\]
we get
\[
f
=
\frac{r^2(\cos^2\theta-\sin^2\theta)}{r^2}
=
\cos2\theta.
\]

The variable \(r\) disappears completely. Therefore the value depends
on the direction \(\theta\).

For \(\theta=0\),
\[
f=1,
\]
while for \(\theta=\frac{\pi}{2}\),
\[
f=-1.
\]

Hence the joint limit does not exist.

This illustrates a useful principle: if polar conversion leaves a
nonconstant angular factor with no vanishing power of \(r\), the limit
is usually path-dependent.

\section{Repeated Limits in Full Detail}

Consider
\[
\lim_{x\to a}
\left[
\lim_{y\to b}f(x,y)
\right].
\]

The first limit treats \(x\) as a fixed parameter. It produces a new
one-variable function
\[
\phi(x)=\lim_{y\to b}f(x,y),
\]
provided this inner limit exists for \(x\) near \(a\).

The outer limit is then
\[
\lim_{x\to a}\phi(x).
\]

The reverse repeated limit may be different because the order of the
two limiting operations has changed.

\begin{example}
Let
\[
f(x,y)=\frac{x-y}{x+y}.
\]

For fixed \(x\ne0\),
\[
\lim_{y\to0}f(x,y)=1.
\]
Hence
\[
\lim_{x\to0}\lim_{y\to0}f(x,y)=1.
\]

For fixed \(y\ne0\),
\[
\lim_{x\to0}f(x,y)=-1.
\]
Hence
\[
\lim_{y\to0}\lim_{x\to0}f(x,y)=-1.
\]

Thus the order matters.
\end{example}

\section{Joint Limit versus Repeated Limits: Logical Summary}

The relationships can be summarized as follows.

\begin{enumerate}
\item If the joint limit exists, then every path limit must agree with
it.

\item If two path limits disagree, the joint limit does not exist.

\item If the two repeated limits disagree, the joint limit cannot
exist.

\item If the two repeated limits agree, the joint limit may still fail
to exist.

\item Agreement along every straight line is still not sufficient for
existence.

\item To establish existence, use a uniform argument that controls all
directions simultaneously.
\end{enumerate}

\section{Continuity on a Region}

A function \(f\) is continuous on a region \(D\) if it is continuous
at every point of \(D\).

For a rational function
\[
f(x,y)=\frac{P(x,y)}{Q(x,y)},
\]
where \(P,Q\) are polynomials, continuity holds at every point where
\[
Q(x,y)\ne0.
\]

For example,
\[
f(x,y)=\frac{x+y}{1+x^2+y^2}
\]
is continuous everywhere because
\[
1+x^2+y^2>0
\]
for all \((x,y)\).

\section{Removable and Nonremovable Discontinuity in Two Variables}

Consider
\[
f(x,y)
=
\frac{x^2+y^2}{\sqrt{x^2+y^2}}
\]
for \((x,y)\ne(0,0)\).

This simplifies to
\[
f(x,y)=\sqrt{x^2+y^2}.
\]

Hence
\[
\lim_{(x,y)\to(0,0)}f(x,y)=0.
\]

If we define \(f(0,0)=0\), the discontinuity is removed.

By contrast,
\[
\frac{x^2-y^2}{x^2+y^2}
\]
has no limit at the origin, so no choice of \(f(0,0)\) can make it
continuous there.

\section{Further Detailed Solved Problems}

\subsection*{Problem 9. Limit by polar coordinates}

Evaluate
\[
\lim_{(x,y)\to(0,0)}
\frac{x^4y^2}{(x^2+y^2)^2}.
\]

Put
\[
x=r\cos\theta,\qquad y=r\sin\theta.
\]

Then
\[
x^4y^2
=
r^6\cos^4\theta\sin^2\theta,
\]
and
\[
(x^2+y^2)^2=r^4.
\]

Therefore
\[
f
=
r^2\cos^4\theta\sin^2\theta.
\]

Since
\[
|\cos^4\theta\sin^2\theta|\le1,
\]
we get
\[
|f|\le r^2.
\]

Thus
\[
\lim_{(x,y)\to(0,0)}
\frac{x^4y^2}{(x^2+y^2)^2}=0.
\]

\subsection*{Problem 10. Detecting a balancing path}

Examine
\[
\lim_{(x,y)\to(0,0)}
\frac{x^3y}{x^6+y^2}.
\]

The denominator contains \(x^6\) and \(y^2\). These have the same order
when
\[
y=kx^3.
\]

Along this path,
\[
\frac{x^3(kx^3)}{x^6+k^2x^6}
=
\frac{k}{1+k^2}.
\]

The result depends on \(k\). Therefore the joint limit does not exist.

\subsection*{Problem 11. A continuous piecewise function}

Let
\[
f(x,y)
=
\begin{cases}
\dfrac{x^2y}{\sqrt{x^2+y^2}},&(x,y)\ne(0,0),\\[1ex]
0,&(x,y)=(0,0).
\end{cases}
\]

For \((x,y)\ne(0,0)\),
\[
|f(x,y)|
=
\frac{|x|^2|y|}{\sqrt{x^2+y^2}}.
\]

Let
\[
r=\sqrt{x^2+y^2}.
\]
Since
\[
|x|\le r,\qquad |y|\le r,
\]
we have
\[
|f(x,y)|
\le
\frac{r^3}{r}
=
r^2.
\]

Hence
\[
f(x,y)\to0=f(0,0).
\]

Therefore \(f\) is continuous at the origin.


\section{Detailed Solved Problems}

\subsection*{Problem 1. A direct limit}
Evaluate
\[
\lim_{(x,y)\to(1,2)}(x^2+xy+y^2).
\]
Because a polynomial is continuous,
\[
1^2+(1)(2)+2^2=7.
\]

\subsection*{Problem 2. Limit by comparison}
Evaluate
\[
\lim_{(x,y)\to(0,0)}
\frac{x^3}{x^2+y^2}.
\]
Since
\[
\abs{\frac{x^3}{x^2+y^2}}
=
\abs{x}\frac{x^2}{x^2+y^2}
\le\abs{x},
\]
and \(\abs{x}\to0\), the limit is \(0\).

\subsection*{Problem 3 (2024). Nonexistence of a double limit}
Examine
\[
\lim_{(x,y)\to(0,0)}
\frac{2x}{x^2+y^2}.
\]
Along \(x=0\), the expression equals \(0\). Along \(y=0\),
\[
\frac{2x}{x^2}=\frac2x,
\]
which has no finite limit as \(x\to0\). Therefore the joint limit
does not exist.

\subsection*{Problem 4. Repeated limits}
Let
\[
f(x,y)=\frac{x-y}{x+y}.
\]
For fixed \(x\ne0\),
\[
\lim_{y\to0}f(x,y)=1,
\]
hence
\[
\lim_{x\to0}\lim_{y\to0}f(x,y)=1.
\]
For fixed \(y\ne0\),
\[
\lim_{x\to0}f(x,y)=-1,
\]
hence
\[
\lim_{y\to0}\lim_{x\to0}f(x,y)=-1.
\]

\section{Graded Practice Problems}

\subsection*{Level I}
\begin{enumerate}
\item Evaluate
\[
\lim_{(x,y)\to(1,1)}(x^2+3xy-y^2).
\]
\item Evaluate
\[
\lim_{(x,y)\to(0,0)}
\frac{x^2y}{x^2+y^2}.
\]
\item Test
\[
\lim_{(x,y)\to(0,0)}
\frac{x^2-y^2}{x^2+y^2}.
\]
\item State the definition of continuity at \((a,b)\).
\end{enumerate}

\subsection*{Level II}
\begin{enumerate}
\item Examine
\[
\lim_{(x,y)\to(0,0)}
\frac{xy}{x^2+y^2}.
\]
\item Prove by the \(\eps\)-\(\delta\) definition that
\[
\lim_{(x,y)\to(0,0)}(x+y)=0.
\]
\item Evaluate both repeated limits of
\[
\frac{x^2-y^2}{x^2+y^2}.
\]
\end{enumerate}

\subsection*{Level III}
\begin{enumerate}
\item Show that straight-line checking is insufficient for
\[
\frac{x^2y}{x^4+y^2}.
\]
\item Find an example where both repeated limits are equal but the
joint limit fails to exist.
\item Use polar coordinates to prove
\[
\lim_{(x,y)\to(0,0)}
\frac{x^4y^2}{x^4+y^2}=0.
\]
\end{enumerate}

% ============================================================
\chapter{Partial Derivatives of First and Second Orders}
% ============================================================

\section*{Learning Objectives}
Students should be able to:
\begin{itemize}
\item define partial derivatives from first principles;
\item interpret them geometrically as slopes of traces;
\item calculate first and second partial derivatives;
\item distinguish pure and mixed derivatives;
\item understand what the existence of partial derivatives does and
does not imply.
\end{itemize}

\section*{Historical Motivation}

A function \(z=f(x,y)\) may respond differently to changes in \(x\)
and changes in \(y\). Partial differentiation isolates one direction
at a time by temporarily fixing the other variable. This became
fundamental in mechanics, thermodynamics and differential equations.

\section{First-Order Partial Derivatives}

\begin{definition}
The partial derivative of \(f\) with respect to \(x\) at \((a,b)\) is
\[
f_x(a,b)
=
\lim_{h\to0}
\frac{f(a+h,b)-f(a,b)}{h},
\]
provided the limit exists.
\end{definition}

Similarly,
\[
f_y(a,b)
=
\lim_{k\to0}
\frac{f(a,b+k)-f(a,b)}{k}.
\]

When differentiating with respect to \(x\), treat \(y\) as a constant.
When differentiating with respect to \(y\), treat \(x\) as a constant.

\paragraph{Real-life comparison.}
If \(T(x,y)\) is the temperature at position \((x,y)\), then \(T_x\)
measures how temperature changes when moving in the \(x\)-direction
while keeping \(y\) fixed; \(T_y\) does the analogous job in the
\(y\)-direction.

\section{Geometrical Interpretation}

For the surface
\[
z=f(x,y),
\]
fix \(y=b\). Then
\[
z=f(x,b)
\]
is a curve in the vertical plane \(y=b\), and
\[
f_x(a,b)
\]
is the tangent slope of this trace at \(x=a\).

Similarly \(f_y(a,b)\) is the tangent slope of the trace obtained by
fixing \(x=a\).


\section{Visualizing the Two Coordinate Slices}

At a point \((a,b)\), the partial derivative \(f_x(a,b)\) is obtained
from the slice in which \(y=b\) is held fixed. Likewise,
\(f_y(a,b)\) comes from the slice \(x=a\).

\begin{center}
\begin{tikzpicture}[scale=1.0]
% left
\begin{scope}[xshift=-3.5cm]
\draw[->] (-2.2,0)--(2.3,0) node[right] {$x$};
\draw[->] (0,-0.4)--(0,3.1) node[above] {$z$};
\draw[thick,domain=-1.8:1.8,samples=50]
plot (\x,{0.55*(\x)^2+0.5});
\fill (0.8,{0.55*0.64+0.5}) circle (2pt);
\draw[thick] (-0.2,0.38)--(1.8,1.5);
\node at (0,-0.9) {\(\text{slice }y=b\)};
\node at (0,-1.35) {\(f_x(a,b)=\text{slope}\)};
\end{scope}

% right
\begin{scope}[xshift=3.5cm]
\draw[->] (-2.2,0)--(2.3,0) node[right] {$y$};
\draw[->] (0,-0.4)--(0,3.1) node[above] {$z$};
\draw[thick,domain=-1.8:1.8,samples=50]
plot (\x,{0.35*(\x)^2+0.7});
\fill (-0.7,{0.35*0.49+0.7}) circle (2pt);
\draw[thick] (-1.6,1.55)--(0.2,0.65);
\node at (0,-0.9) {\(\text{slice }x=a\)};
\node at (0,-1.35) {\(f_y(a,b)=\text{slope}\)};
\end{scope}
\end{tikzpicture}
\end{center}

Each partial derivative is therefore an ordinary one-variable tangent
slope on a particular vertical cross-section of the surface.


\section{Calculation from the Definition}

Let
\[
f(x,y)=x^2+y^2.
\]
Then
\[
\begin{aligned}
f_x(a,b)
&=
\lim_{h\to0}
\frac{(a+h)^2+b^2-(a^2+b^2)}{h}\\
&=
\lim_{h\to0}
\frac{2ah+h^2}{h}\\
&=
\lim_{h\to0}(2a+h)\\
&=2a.
\end{aligned}
\]
Similarly,
\[
f_y(a,b)=2b.
\]

\section{Second-Order Partial Derivatives}

The four second-order partial derivatives are
\[
f_{xx},\qquad f_{xy},\qquad f_{yx},\qquad f_{yy}.
\]

Here
\[
f_{xy}
=
\frac{\partial}{\partial y}(f_x),
\qquad
f_{yx}
=
\frac{\partial}{\partial x}(f_y).
\]

\begin{example}
Let
\[
f(x,y)=x^3y^2+e^{xy}.
\]
Then
\[
f_x=3x^2y^2+ye^{xy},
\]
\[
f_y=2x^3y+xe^{xy}.
\]
Differentiate again:
\[
f_{xx}=6xy^2+y^2e^{xy},
\]
\[
f_{xy}=6x^2y+e^{xy}+xye^{xy},
\]
\[
f_{yx}=6x^2y+e^{xy}+xye^{xy},
\]
and
\[
f_{yy}=2x^3+x^2e^{xy}.
\]
Thus in this example
\[
f_{xy}=f_{yx}.
\]
\end{example}

\section{Partial Derivatives Need Not Imply Continuity}

Define
\[
f(x,y)=
\begin{cases}
\dfrac{xy}{x^2+y^2},&(x,y)\ne(0,0),\\[1ex]
0,&(x,y)=(0,0).
\end{cases}
\]

By definition,
\[
f_x(0,0)
=
\lim_{h\to0}
\frac{f(h,0)-f(0,0)}{h}
=0.
\]
Likewise,
\[
f_y(0,0)=0.
\]

But along \(y=x\),
\[
f(x,x)=\frac12,
\]
so the function is not continuous at \((0,0)\).

Therefore
\[
f_x(0,0),f_y(0,0)\text{ exist}
\quad\not\Rightarrow\quad
f\text{ is continuous at }(0,0).
\]


\section{Notation for Partial Derivatives}

Several notations are used interchangeably:
\[
f_x
=
\pd{f}{x},
\qquad
f_y
=
\pd{f}{y}.
\]

For second derivatives,
\[
f_{xx}
=
\frac{\partial^2f}{\partial x^2},
\qquad
f_{yy}
=
\frac{\partial^2f}{\partial y^2},
\]
while
\[
f_{xy}
=
\frac{\partial}{\partial y}
\left(
\pd{f}{x}
\right),
\]
and
\[
f_{yx}
=
\frac{\partial}{\partial x}
\left(
\pd{f}{y}
\right).
\]

Students should pay attention to the order: \(f_{xy}\) means
differentiate first with respect to \(x\), then with respect to \(y\).

\section{Partial Derivatives as Rates of Change}

Suppose
\[
T=T(x,y)
\]
describes temperature on a thin metal plate.

Then
\[
T_x(a,b)
\]
is the instantaneous rate of change of temperature when we move in the
\(x\)-direction through \((a,b)\), keeping \(y=b\) fixed.

Similarly,
\[
T_y(a,b)
\]
describes the rate of change in the \(y\)-direction while keeping
\(x=a\) fixed.

Thus partial derivatives are ordinary one-variable derivatives of
appropriate slices of the surface.

\section{A Detailed First-Principles Example}

Let
\[
f(x,y)=xy+x^2.
\]
Find \(f_x(a,b)\) directly from the definition.

By definition,
\[
f_x(a,b)
=
\lim_{h\to0}
\frac{
f(a+h,b)-f(a,b)
}{h}.
\]

Now
\[
f(a+h,b)
=
(a+h)b+(a+h)^2.
\]
Therefore
\[
\begin{aligned}
f(a+h,b)-f(a,b)
&=
ab+hb+a^2+2ah+h^2-ab-a^2\\
&=
hb+2ah+h^2.
\end{aligned}
\]

Thus
\[
\frac{
f(a+h,b)-f(a,b)
}{h}
=
b+2a+h.
\]

Letting \(h\to0\),
\[
f_x(a,b)=b+2a.
\]

The same result is obtained by ordinary partial differentiation:
\[
f_x=y+2x.
\]

\section{Higher-Order Derivatives and Harmonic Functions}

Second derivatives describe how the first rates of change themselves
vary.

A particularly important combination is the two-dimensional Laplacian
\[
\nabla^2u
=
u_{xx}+u_{yy}.
\]

\begin{definition}
A twice continuously differentiable function \(u(x,y)\) is called
\emph{harmonic} in a region if
\[
u_{xx}+u_{yy}=0
\]
throughout that region.
\end{definition}

\begin{example}
Let
\[
u(x,y)=e^x\cos y.
\]
Then
\[
u_x=e^x\cos y,
\qquad
u_{xx}=e^x\cos y.
\]
Also,
\[
u_y=-e^x\sin y,
\qquad
u_{yy}=-e^x\cos y.
\]
Therefore
\[
u_{xx}+u_{yy}=0.
\]
Hence \(u\) is harmonic.
\end{example}

\section{Additional Detailed Solved Problems}

\subsection*{Problem 3. First and second partial derivatives}

Let
\[
u=x^4y^3+\sin(xy).
\]

First,
\[
u_x
=
4x^3y^3+y\cos(xy),
\]
and
\[
u_y
=
3x^4y^2+x\cos(xy).
\]

Now
\[
u_{xx}
=
12x^2y^3-y^2\sin(xy),
\]
and
\[
u_{yy}
=
6x^4y-x^2\sin(xy).
\]

For the mixed derivative,
\[
\begin{aligned}
u_{xy}
&=
\frac{\partial}{\partial y}
\left(
4x^3y^3+y\cos(xy)
\right)\\
&=
12x^3y^2+\cos(xy)-xy\sin(xy).
\end{aligned}
\]

Similarly,
\[
\begin{aligned}
u_{yx}
&=
\frac{\partial}{\partial x}
\left(
3x^4y^2+x\cos(xy)
\right)\\
&=
12x^3y^2+\cos(xy)-xy\sin(xy).
\end{aligned}
\]

Hence
\[
u_{xy}=u_{yx}.
\]

\subsection*{Problem 4. Partial derivatives exist but continuity fails}

Consider
\[
f(x,y)
=
\begin{cases}
\dfrac{xy}{x^2+y^2},&(x,y)\ne(0,0),\\
0,&(x,y)=(0,0).
\end{cases}
\]

At the origin,
\[
f_x(0,0)
=
\lim_{h\to0}
\frac{f(h,0)-f(0,0)}{h}
=0,
\]
because \(f(h,0)=0\).

Likewise,
\[
f_y(0,0)=0.
\]

But along \(y=x\),
\[
f(x,x)=\frac12.
\]

Therefore the function is not continuous at the origin. This example
shows clearly that the mere existence of both first partial derivatives
is much weaker than differentiability.



\section{Partial Derivatives as Ordinary Derivatives of Slices}

Suppose
\[
z=f(x,y).
\]

To compute \(f_x(a,b)\), hold \(y=b\) fixed and define the one-variable
function
\[
g(x)=f(x,b).
\]

Then
\[
f_x(a,b)=g'(a).
\]

Similarly, if
\[
h(y)=f(a,y),
\]
then
\[
f_y(a,b)=h'(b).
\]

This viewpoint is useful because partial differentiation is not a new
kind of differentiation. It is ordinary one-variable differentiation
applied to a suitable cross-section of the surface.

\section{Partial Derivatives from First Principles: A Second Example}

Let
\[
f(x,y)=x^2y+xy^2.
\]

Find \(f_y(a,b)\) directly from the definition.

By definition,
\[
f_y(a,b)
=
\lim_{k\to0}
\frac{
f(a,b+k)-f(a,b)
}{k}.
\]

Now
\[
f(a,b+k)
=
a^2(b+k)+a(b+k)^2.
\]

Expanding,
\[
\begin{aligned}
f(a,b+k)
&=
a^2b+a^2k+a(b^2+2bk+k^2)\\
&=
a^2b+a^2k+ab^2+2abk+ak^2.
\end{aligned}
\]

Since
\[
f(a,b)=a^2b+ab^2,
\]
we obtain
\[
f(a,b+k)-f(a,b)
=
a^2k+2abk+ak^2.
\]

Divide by \(k\):
\[
\frac{
f(a,b+k)-f(a,b)
}{k}
=
a^2+2ab+ak.
\]

Letting \(k\to0\),
\[
f_y(a,b)=a^2+2ab.
\]

Thus
\[
f_y=x^2+2xy,
\]
as obtained by the usual rule.

\section{First-Order Increment Approximation}

If \(f\) is differentiable, then for small changes
\[
\Delta x=h,\qquad \Delta y=k,
\]
the corresponding change in \(f\) is approximately
\[
\Delta f
\approx
f_x\,\Delta x+f_y\,\Delta y.
\]

Even before formal differentiability is introduced, this formula
provides intuition: \(f_x\) measures the contribution from the
\(x\)-change, while \(f_y\) measures the contribution from the
\(y\)-change.

\begin{example}
Let
\[
f(x,y)=x^2y.
\]
At \((1,2)\),
\[
f_x=2xy,
\qquad
f_y=x^2.
\]

Hence
\[
f_x(1,2)=4,
\qquad
f_y(1,2)=1.
\]

If
\[
\Delta x=0.01,\qquad
\Delta y=-0.02,
\]
then the first-order prediction is
\[
\Delta f
\approx
4(0.01)+1(-0.02)
=
0.02.
\]
\end{example}

\section{Second Partial Derivatives: Order Matters in the Notation}

The mixed derivatives are
\[
f_{xy}
=
\frac{\partial}{\partial y}
\left(
\frac{\partial f}{\partial x}
\right),
\]
and
\[
f_{yx}
=
\frac{\partial}{\partial x}
\left(
\frac{\partial f}{\partial y}
\right).
\]

Thus the order of the subscripts records the order of differentiation.

For sufficiently smooth functions these are equal, but this equality
is a theorem, not a consequence of notation.

\section{A Detailed Mixed-Derivative Computation}

Let
\[
f(x,y)=x^2e^{xy}.
\]

First,
\[
f_x
=
2xe^{xy}+x^2ye^{xy}.
\]

Now differentiate with respect to \(y\):
\[
\begin{aligned}
f_{xy}
&=
2x(xe^{xy})
+
x^2\frac{\partial}{\partial y}
\left(ye^{xy}\right)\\
&=
2x^2e^{xy}
+
x^2\left(e^{xy}+xy e^{xy}\right)\\
&=
3x^2e^{xy}
+
x^3y e^{xy}.
\end{aligned}
\]

Now start in the opposite order.

\[
f_y
=
x^3e^{xy}.
\]

Hence
\[
\begin{aligned}
f_{yx}
&=
3x^2e^{xy}
+
x^3ye^{xy}.
\end{aligned}
\]

Therefore
\[
f_{xy}=f_{yx}.
\]

\section{Common Errors in Partial Differentiation}

\begin{enumerate}
\item Forgetting which variable is held constant.

For example, if
\[
f(x,y)=x^2y^3,
\]
then
\[
f_x=2xy^3,
\]
not
\[
2xy^2.
\]

\item Treating \(xy\) as a constant when differentiating
\(e^{xy}\).

For \(x\)-differentiation,
\[
\frac{\partial}{\partial x}e^{xy}
=
ye^{xy}.
\]

\item Forgetting product or chain rules in partial derivatives.

For
\[
f(x,y)=x^2\sin(xy),
\]
\[
f_x
=
2x\sin(xy)+x^2y\cos(xy).
\]

\item Assuming automatically that
\[
f_{xy}=f_{yx}.
\]
Equality requires appropriate hypotheses.
\end{enumerate}

\section{More Detailed Solved Problems}

\subsection*{Problem 5. Compute all second-order partial derivatives}

Let
\[
u=\log(x^2+y^2),
\]
where \((x,y)\ne(0,0)\).

First,
\[
u_x
=
\frac{2x}{x^2+y^2},
\qquad
u_y
=
\frac{2y}{x^2+y^2}.
\]

Differentiate again:
\[
u_{xx}
=
\frac{
2(x^2+y^2)-4x^2
}{
(x^2+y^2)^2
}
=
\frac{2(y^2-x^2)}{(x^2+y^2)^2}.
\]

Similarly,
\[
u_{yy}
=
\frac{2(x^2-y^2)}{(x^2+y^2)^2}.
\]

For the mixed derivative,
\[
u_{xy}
=
-\frac{4xy}{(x^2+y^2)^2}.
\]

Likewise,
\[
u_{yx}
=
-\frac{4xy}{(x^2+y^2)^2}.
\]

Thus
\[
u_{xy}=u_{yx}.
\]

Also,
\[
u_{xx}+u_{yy}=0.
\]

Hence \(\log(x^2+y^2)\) is harmonic away from the origin.

\subsection*{Problem 6. First-principles partial derivative of a piecewise function}

Let
\[
f(x,y)
=
\begin{cases}
\dfrac{x^2y}{x^2+y^2},&(x,y)\ne(0,0),\\
0,&(x,y)=(0,0).
\end{cases}
\]

Find \(f_x(0,0)\).

By definition,
\[
f_x(0,0)
=
\lim_{h\to0}
\frac{f(h,0)-f(0,0)}{h}.
\]

Since
\[
f(h,0)=0,
\]
we get
\[
f_x(0,0)=0.
\]

Similarly,
\[
f_y(0,0)=0.
\]

This calculation illustrates why coordinate partial derivatives can
exist even when the full two-variable behavior is more complicated.


\section{Detailed Solved Problems}

\subsection*{Problem 1}
For
\[
u=x^2y^3+e^{x+y},
\]
find all first and second partial derivatives.

First:
\[
u_x=2xy^3+e^{x+y},
\qquad
u_y=3x^2y^2+e^{x+y}.
\]
Second:
\[
u_{xx}=2y^3+e^{x+y},
\]
\[
u_{xy}=6xy^2+e^{x+y},
\]
\[
u_{yx}=6xy^2+e^{x+y},
\]
\[
u_{yy}=6x^2y+e^{x+y}.
\]

\subsection*{Problem 2}
Find \(f_x(0,0)\) and \(f_y(0,0)\) directly for
\[
f(x,y)=
\begin{cases}
\dfrac{x^2y}{x^4+y^2},&(x,y)\ne(0,0),\\
0,&(0,0).
\end{cases}
\]
Since \(f(h,0)=0\),
\[
f_x(0,0)
=
\lim_{h\to0}\frac{0-0}{h}=0.
\]
Similarly \(f(0,k)=0\), so
\[
f_y(0,0)=0.
\]

\section{Graded Practice Problems}
\begin{enumerate}
\item Find all first partial derivatives of
\[
x^3+3x^2y+y^3.
\]
\item Find all second partial derivatives of \(e^{x+y}\).
\item Find \(f_x(1,2)\) and \(f_y(1,2)\) for \(x^2y+xy^2\).
\item Compute \(f_x(0,0),f_y(0,0)\) from first principles for
\(x^2+y^2\).
\item Give an example where both first partial derivatives exist but
the function is not continuous.
\end{enumerate}

% ============================================================
\chapter{Differentiability, Young's Theorem and Schwarz's Theorem}
% ============================================================

\section*{Learning Objectives}
Students should be able to:
\begin{itemize}
\item state the definition of differentiability in two variables;
\item interpret differentiability as linear approximation;
\item prove that differentiability implies continuity;
\item use continuity of first partial derivatives as a sufficient
condition for differentiability;
\item state Young's and Schwarz's theorems;
\item apply equality of mixed partial derivatives correctly.
\end{itemize}

\section*{Why Partial Derivatives Are Not Enough}

Partial derivatives only test changes along the coordinate directions.
Differentiability requires one linear approximation that works for
\emph{all} sufficiently small changes \((h,k)\).

\section{Definition of Differentiability}

\begin{definition}
A function \(f(x,y)\) is differentiable at \((a,b)\) if
\[
f(a+h,b+k)-f(a,b)
=
f_x(a,b)h+f_y(a,b)k
+
o\!\left(\sqrt{h^2+k^2}\right).
\]
Equivalently,
\[
\lim_{(h,k)\to(0,0)}
\frac{
f(a+h,b+k)-f(a,b)
-f_x(a,b)h-f_y(a,b)k
}{
\sqrt{h^2+k^2}
}
=0.
\]
\end{definition}

\paragraph{Geometrical meaning.}
Near \((a,b)\),
\[
f(a+h,b+k)
\approx
f(a,b)+f_x(a,b)h+f_y(a,b)k.
\]
Thus the surface is locally approximated by a plane.

The tangent plane is
\[
z
=
f(a,b)
+
f_x(a,b)(x-a)
+
f_y(a,b)(y-b).
\]

\section{Differentiability Implies Partial Derivatives Exist}

\begin{theorem}
If \(f\) is differentiable at \((a,b)\), then both \(f_x(a,b)\) and
\(f_y(a,b)\) exist.
\end{theorem}

\begin{proof}
Suppose
\[
f(a+h,b+k)-f(a,b)
=
Ah+Bk+\eta(h,k)\sqrt{h^2+k^2},
\]
where
\[
\eta(h,k)\to0.
\]

Set \(k=0\):
\[
f(a+h,b)-f(a,b)
=
Ah+\eta(h,0)\abs{h}.
\]
For \(h\ne0\),
\[
\frac{f(a+h,b)-f(a,b)}{h}
=
A+\eta(h,0)\frac{\abs h}{h}.
\]
Since
\[
\left|
\eta(h,0)\frac{\abs h}{h}
\right|
=
\abs{\eta(h,0)}
\to0,
\]
we obtain
\[
f_x(a,b)=A.
\]

Setting \(h=0\) similarly gives
\[
f_y(a,b)=B.
\]
\end{proof}

\section{Every Differentiable Function Is Continuous}

\begin{theorem}
If \(f(x,y)\) is differentiable at \((a,b)\), then \(f\) is
continuous at \((a,b)\).
\end{theorem}

\begin{proof}
To prove continuity we must show
\[
\lim_{(x,y)\to(a,b)}f(x,y)=f(a,b).
\]

Put
\[
x=a+h,\qquad y=b+k.
\]
Then
\[
(x,y)\to(a,b)
\quad\Longleftrightarrow\quad
(h,k)\to(0,0).
\]

Since \(f\) is differentiable,
\[
f(a+h,b+k)-f(a,b)
=
f_x(a,b)h
+
f_y(a,b)k
+
\eta(h,k)\sqrt{h^2+k^2},
\]
where
\[
\eta(h,k)\to0.
\]

Taking absolute values,
\[
\begin{aligned}
&
\abs{f(a+h,b+k)-f(a,b)}
\\
&\le
\abs{f_x(a,b)}\abs h
+
\abs{f_y(a,b)}\abs k
+
\abs{\eta(h,k)}\sqrt{h^2+k^2}.
\end{aligned}
\]

Since
\[
\abs h,\abs k\le\sqrt{h^2+k^2},
\]
we get
\[
\begin{aligned}
&
\abs{f(a+h,b+k)-f(a,b)}
\\
&\le
\left(
\abs{f_x(a,b)}
+\abs{f_y(a,b)}
+\abs{\eta(h,k)}
\right)
\sqrt{h^2+k^2}.
\end{aligned}
\]

As \((h,k)\to(0,0)\), the right-hand side tends to \(0\).
Therefore
\[
f(a+h,b+k)-f(a,b)\to0.
\]
Hence
\[
\lim_{(h,k)\to(0,0)}f(a+h,b+k)=f(a,b).
\]
Equivalently,
\[
\lim_{(x,y)\to(a,b)}f(x,y)=f(a,b).
\]
Thus \(f\) is continuous at \((a,b)\).
\end{proof}

\section{A Standard Sufficient Condition}

\begin{theorem}
Suppose \(f_x\) and \(f_y\) exist in a neighbourhood of \((a,b)\)
and are continuous at \((a,b)\). Then \(f\) is differentiable at
\((a,b)\).
\end{theorem}

\begin{proof}
Write
\[
\Delta f=f(a+h,b+k)-f(a,b).
\]
Insert and subtract \(f(a,b+k)\):
\[
\begin{aligned}
\Delta f
={}&
[f(a+h,b+k)-f(a,b+k)]\\
&+[f(a,b+k)-f(a,b)].
\end{aligned}
\]

Apply the one-variable Mean Value Theorem to the first difference as a
function of \(x\). For some \(\theta\in(0,1)\),
\[
f(a+h,b+k)-f(a,b+k)
=
h f_x(a+\theta h,b+k).
\]

Apply the Mean Value Theorem to the second difference as a function of
\(y\). For some \(\lambda\in(0,1)\),
\[
f(a,b+k)-f(a,b)
=
k f_y(a,b+\lambda k).
\]

Therefore
\[
\Delta f
=
h f_x(a+\theta h,b+k)
+
k f_y(a,b+\lambda k).
\]

Subtract
\[
f_x(a,b)h+f_y(a,b)k.
\]
Then
\[
\begin{aligned}
&
\Delta f-f_x(a,b)h-f_y(a,b)k\\
={}&
h\bigl[
f_x(a+\theta h,b+k)-f_x(a,b)
\bigr]\\
&+
k\bigl[
f_y(a,b+\lambda k)-f_y(a,b)
\bigr].
\end{aligned}
\]

Let
\[
\rho=\sqrt{h^2+k^2}.
\]
Then
\[
\frac{\abs h}{\rho}\le1,
\qquad
\frac{\abs k}{\rho}\le1.
\]
Hence the absolute value of the differentiability quotient is bounded
by
\[
\abs{f_x(a+\theta h,b+k)-f_x(a,b)}
+
\abs{f_y(a,b+\lambda k)-f_y(a,b)}.
\]
By continuity of \(f_x,f_y\), this tends to \(0\). Therefore \(f\) is
differentiable at \((a,b)\).
\end{proof}

\section{Logical Relationships}

It is useful to remember:
\[
\text{continuous first partial derivatives near the point}
\Longrightarrow
\text{differentiable}
\]
and
\[
\text{differentiable}
\Longrightarrow
\text{continuous}.
\]

But
\[
\text{partial derivatives exist}
\not\Rightarrow
\text{continuous}
\]
and therefore do not imply differentiability.


\section{A Stronger Counterexample: Continuity and Partial Derivatives Still Need Not Give Differentiability}

The following example is particularly important because it shows that
even the combination
\[
\text{continuity at the point}
\quad+\quad
\text{existence of both first partial derivatives}
\]
is still not enough to guarantee differentiability.

Define
\[
f(x,y)=
\begin{cases}
\dfrac{x^3}{x^2+y^2},&(x,y)\ne(0,0),\\[1ex]
0,&(x,y)=(0,0).
\end{cases}
\]

We examine the function carefully.

\subsection*{Step 1. Continuity at the origin}

For \((x,y)\ne(0,0)\),
\[
\left|
\frac{x^3}{x^2+y^2}
\right|
=
|x|\frac{x^2}{x^2+y^2}.
\]

Since
\[
0\le
\frac{x^2}{x^2+y^2}
\le1,
\]
we obtain
\[
|f(x,y)|\le |x|.
\]

As
\[
(x,y)\to(0,0),
\]
we have
\[
|x|\to0.
\]

Hence, by the Sandwich Theorem,
\[
\lim_{(x,y)\to(0,0)}f(x,y)=0=f(0,0).
\]

Therefore \(f\) is continuous at \((0,0)\).

\subsection*{Step 2. The partial derivative \(f_x(0,0)\)}

By definition,
\[
f_x(0,0)
=
\lim_{h\to0}
\frac{f(h,0)-f(0,0)}{h}.
\]

Now
\[
f(h,0)
=
\frac{h^3}{h^2}
=
h
\qquad(h\ne0).
\]

Therefore
\[
f_x(0,0)
=
\lim_{h\to0}\frac{h}{h}
=
1.
\]

\subsection*{Step 3. The partial derivative \(f_y(0,0)\)}

Similarly,
\[
f_y(0,0)
=
\lim_{k\to0}
\frac{f(0,k)-f(0,0)}{k}.
\]

But
\[
f(0,k)=0.
\]

Hence
\[
f_y(0,0)=0.
\]

Thus both first partial derivatives exist.

\subsection*{Step 4. Test differentiability}

If \(f\) were differentiable at the origin, then we would need
\[
\frac{
f(h,k)-f(0,0)-f_x(0,0)h-f_y(0,0)k
}{
\sqrt{h^2+k^2}
}
\longrightarrow0.
\]

Since
\[
f(0,0)=0,\qquad
f_x(0,0)=1,\qquad
f_y(0,0)=0,
\]
this becomes
\[
\frac{
f(h,k)-h
}{
\sqrt{h^2+k^2}}
\longrightarrow0.
\]

Now choose the path
\[
k=h.
\]

Then
\[
f(h,h)
=
\frac{h^3}{2h^2}
=
\frac h2.
\]

Therefore
\[
\frac{
f(h,h)-h
}{
\sqrt{h^2+h^2}}
=
\frac{-h/2}{\sqrt2\,|h|}.
\]

Its absolute value is
\[
\frac{1}{2\sqrt2},
\]
which does not tend to \(0\).

Therefore \(f\) is \emph{not differentiable} at the origin.

So this example proves
\[
\text{continuous at the point + partial derivatives exist}
\not\Longrightarrow
\text{differentiable}.
\]

\section{Logical Map of the Main Concepts}

The following diagram summarizes the implications that are valid under
the hypotheses used in this chapter.

\begin{center}
\begin{tikzpicture}[
node distance=1.7cm and 2.0cm,
every node/.style={align=center},
arr/.style={-{Stealth[length=2.5mm]},thick}
]
\node (cp) {\textbf{First partial derivatives}\\
exist in a neighbourhood\\
and are continuous at the point};

\node[below=of cp] (d) {\textbf{Differentiable}\\
at the point};

\node[below left=1.7cm and 1.2cm of d] (c) {\textbf{Continuous}\\
at the point};

\node[below right=1.7cm and 1.2cm of d] (p) {\textbf{First partial derivatives}\\
exist at the point};

\draw[arr] (cp)--(d);
\draw[arr] (d)--(c);
\draw[arr] (d)--(p);

\node[below=1.1cm of c,align=left] (n1)
{\small Converse need not hold.};
\node[below=1.1cm of p,align=left] (n2)
{\small Converse need not hold.};
\end{tikzpicture}
\end{center}

In particular,
\[
\text{differentiable}
\Longrightarrow
\text{continuous},
\]
but
\[
\text{continuous}
\not\Longrightarrow
\text{differentiable}.
\]

Also,
\[
\text{differentiable}
\Longrightarrow
\text{partial derivatives exist},
\]
but the existence of partial derivatives alone does not guarantee
differentiability.

The counterexample above shows something stronger:
even continuity together with existence of both partial derivatives is
still insufficient in general.



\section{Young's Theorem and Schwarz's Theorem}

\subsection*{A note on terminology}

The names \emph{Clairaut's theorem}, \emph{Schwarz's theorem}, and
\emph{Young's theorem} are not used in exactly the same way in every
calculus or analysis textbook. In some elementary texts, ``Clairaut's
theorem'' or ``Schwarz's theorem'' refers simply to the familiar
\(C^2\)-result that continuous mixed second partial derivatives are
equal.

To avoid ambiguity, throughout these notes we shall use the following
precise formulations.

\subsection*{Notation convention}

We use
\[
f_{xy}
=
\frac{\partial}{\partial y}
\left(
\frac{\partial f}{\partial x}
\right),
\qquad
f_{yx}
=
\frac{\partial}{\partial x}
\left(
\frac{\partial f}{\partial y}
\right).
\]

Thus in \(f_{xy}\), differentiation with respect to \(x\) is performed
first and differentiation with respect to \(y\) second.

\section{Young's Theorem}

\begin{theorem}[Young's theorem --- statement]
Let \(f\) be defined in a neighbourhood of \((a,b)\), and suppose that
the first partial derivatives
\[
f_x,\qquad f_y
\]
exist in that neighbourhood. If the functions \(f_x\) and \(f_y\)
are themselves totally differentiable at \((a,b)\), then the two
mixed second partial derivatives at the point are equal:
\[
f_{xy}(a,b)=f_{yx}(a,b).
\]
\end{theorem}

Only the statement and applications are required for the present
syllabus.

\paragraph{What the hypothesis means.}
Young's theorem places the regularity assumption on the
\emph{first-derivative functions} \(f_x\) and \(f_y\). They must behave
linearly, to first order, near \((a,b)\).

In particular, total differentiability of \(f_x\) and \(f_y\) is
stronger than merely saying that the numerical values
\(f_{xy}(a,b)\) and \(f_{yx}(a,b)\) happen to exist.

\section{Schwarz's Theorem}

\begin{theorem}[Schwarz's theorem --- statement]
Let \(f\) be defined in a neighbourhood \(U\) of \((a,b)\). Suppose
that
\[
f_x,\qquad f_y,\qquad f_{xy}
\]
exist throughout \(U\), and suppose that \(f_{xy}\) is continuous at
\((a,b)\). Then \(f_{yx}(a,b)\) exists and
\[
f_{xy}(a,b)=f_{yx}(a,b).
\]
\end{theorem}

Notice an important feature of this statement: we do not need to assume
in advance that \(f_{yx}(a,b)\) exists. Under the hypotheses of the
theorem, its existence follows.

\section{The Common Schwarz--Clairaut Corollary}

For undergraduate calculations, the following more symmetric form is
often the one used most frequently.

\begin{corollary}[Schwarz--Clairaut]
If the mixed second partial derivatives
\[
f_{xy}
\quad\text{and}\quad
f_{yx}
\]
exist in a neighbourhood of \((a,b)\) and are continuous at
\((a,b)\), then
\[
f_{xy}(a,b)=f_{yx}(a,b).
\]
\end{corollary}

A still more convenient sufficient condition is:
\[
f\in C^2
\quad\Longrightarrow\quad
f_{xy}=f_{yx}.
\]

Here \(f\in C^2\) means that all second-order partial derivatives exist
and are continuous.

\section{Young versus Schwarz: What Is the Difference?}

Both theorems lead to the same conclusion,
\[
f_{xy}(a,b)=f_{yx}(a,b),
\]
but they reach it from different regularity assumptions.

\begin{center}
\renewcommand{\arraystretch}{1.35}
\begin{tabularx}{0.94\textwidth}{
>{\bfseries}p{0.20\textwidth}
p{0.36\textwidth}
X}
\toprule
Theorem & Main hypothesis & Conclusion\\
\midrule
Young
&
\(f_x\) and \(f_y\) exist near the point and are totally
differentiable at \((a,b)\)
&
\(
f_{xy}(a,b)=f_{yx}(a,b)
\)
\\
Schwarz
&
\(f_x,f_y,f_{xy}\) exist near the point and \(f_{xy}\) is
continuous at \((a,b)\)
&
\(f_{yx}(a,b)\) exists and
\(
f_{xy}(a,b)=f_{yx}(a,b)
\)
\\
Schwarz--Clairaut corollary
&
Both mixed second partial derivatives are continuous near the point
&
\(
f_{xy}(a,b)=f_{yx}(a,b)
\)
\\
\bottomrule
\end{tabularx}
\end{center}

Thus it is misleading to summarize the distinction merely as
\[
\text{``one mixed partial continuous''}
\quad\text{versus}\quad
\text{``both mixed partials continuous.''}
\]

The more meaningful distinction is that Young's theorem uses
differentiability of the first-derivative functions, whereas the
Schwarz-type results use continuity assumptions on mixed second
partial derivatives.

\section{Application 1: A \(C^2\) Function}

Let
\[
f(x,y)=x^3y^2+e^{xy}.
\]

First,
\[
f_x=3x^2y^2+ye^{xy}.
\]
Therefore
\[
f_{xy}=6x^2y+e^{xy}+xye^{xy}.
\]

Also,
\[
f_y=2x^3y+xe^{xy},
\]
so
\[
f_{yx}=6x^2y+e^{xy}+xye^{xy}.
\]

All second partial derivatives involved are continuous everywhere.
Therefore the Schwarz--Clairaut theorem applies and
\[
f_{xy}=f_{yx}.
\]

\section{Application 2: Schwarz Can Supply Existence}

Suppose \(f_x,f_y,f_{xy}\) exist throughout a neighbourhood of
\((a,b)\), and suppose \(f_{xy}\) is continuous at \((a,b)\).

Then Schwarz's theorem tells us that
\[
f_{yx}(a,b)
\]
exists and
\[
f_{yx}(a,b)=f_{xy}(a,b).
\]

Thus Schwarz's theorem can do more than merely compare two mixed
partials already known to exist.

\section{Why the Hypotheses Cannot Be Omitted}

Consider
\[
f(x,y)=
\begin{cases}
\dfrac{xy(x^2-y^2)}{x^2+y^2},&(x,y)\ne(0,0),\\[1ex]
0,&(x,y)=(0,0).
\end{cases}
\]

For \(y\ne0\),
\[
f_x(0,y)=-y,
\qquad
f_x(0,0)=0.
\]

Therefore
\[
f_{xy}(0,0)
=
\lim_{y\to0}
\frac{-y}{y}
=-1.
\]

Similarly,
\[
f_y(x,0)=x,
\qquad
f_y(0,0)=0,
\]
and so
\[
f_{yx}(0,0)
=
\lim_{x\to0}
\frac{x}{x}
=1.
\]

Hence
\[
f_{xy}(0,0)=-1,
\qquad
f_{yx}(0,0)=1.
\]

Thus mixed partial derivatives can both exist at a point and still be
unequal if the required regularity hypotheses are absent.

\section{How to Use These Theorems in Problems}

When a problem asks whether mixed partial derivatives may be
interchanged, use the following checklist.

\begin{enumerate}
\item If the function is clearly \(C^2\), use Schwarz--Clairaut.

\item If \(f_x,f_y,f_{xy}\) exist near the point and \(f_{xy}\) is
continuous at the point, Schwarz's theorem applies.

\item If \(f_x\) and \(f_y\) are totally differentiable at the point,
Young's theorem applies.

\item If these hypotheses are not available, do not interchange the
order automatically. Compute directly if necessary.
\end{enumerate}


\section{Increment Form of Differentiability}

Let
\[
x=a+h,\qquad y=b+k.
\]
The change in the function is
\[
\Delta f
=
f(a+h,b+k)-f(a,b).
\]

If \(f\) is differentiable at \((a,b)\), then
\[
\Delta f
=
f_x(a,b)h
+
f_y(a,b)k
+
R(h,k),
\]
where the remainder satisfies
\[
\frac{R(h,k)}{\sqrt{h^2+k^2}}
\to0.
\]

Thus the first-order part of the change is
\[
df
=
f_x\,dx+f_y\,dy.
\]

The quantity \(df\) is called the \emph{total differential}. For small
changes,
\[
\Delta f\approx df.
\]

\paragraph{Real-life comparison.}
Suppose the cost \(C(x,y)\) of production depends on two small input
changes \(h\) and \(k\). Differentiability says that, sufficiently near
the operating point, the total change in cost is well approximated by
adding the separate first-order effects
\[
C_xh+C_yk.
\]
The remaining error is much smaller than the size
\[
\sqrt{h^2+k^2}
\]
of the input change.

\section{Tangent Plane and Linear Approximation}

For a differentiable function \(z=f(x,y)\), the tangent plane at
\[
(a,b,f(a,b))
\]
is
\[
z-f(a,b)
=
f_x(a,b)(x-a)
+
f_y(a,b)(y-b).
\]

Equivalently,
\[
z
=
f(a,b)
+
f_x(a,b)(x-a)
+
f_y(a,b)(y-b).
\]

Therefore the linear approximation is
\[
L(x,y)
=
f(a,b)
+
f_x(a,b)(x-a)
+
f_y(a,b)(y-b).
\]

\begin{example}
Find the tangent plane to
\[
z=x^2+xy+y^2
\]
at \((1,1)\).

We have
\[
f(1,1)=3.
\]
Also,
\[
f_x=2x+y,
\]
so
\[
f_x(1,1)=3.
\]
And
\[
f_y=x+2y,
\]
so
\[
f_y(1,1)=3.
\]

Hence
\[
z-3
=
3(x-1)+3(y-1).
\]

Thus
\[
z=3x+3y-3.
\]
\end{example}

\section{Differentiability versus Continuity versus Partial Derivatives}

The correct implication chain is
\[
\text{continuous first partials near the point}
\Longrightarrow
\text{differentiable}
\Longrightarrow
\text{continuous}.
\]

Also,
\[
\text{differentiable}
\Longrightarrow
\text{first partial derivatives exist}.
\]

But none of the reverse implications is automatic.

In particular:

\[
f_x,f_y\text{ exist}
\not\Rightarrow
f\text{ continuous},
\]

and therefore

\[
f_x,f_y\text{ exist}
\not\Rightarrow
f\text{ differentiable}.
\]

Also,
\[
f\text{ continuous}
\not\Rightarrow
f_x,f_y\text{ exist}.
\]

For example,
\[
f(x,y)=\sqrt{x^2+y^2}
\]
is continuous at the origin, but
\[
f_x(0,0)
=
\lim_{h\to0}
\frac{|h|}{h}
\]
does not exist.


\section{Surface and Tangent Plane}

For
\[
f(x,y)=x^2+y^2,
\]
consider the point \((1,1,2)\). Since
\[
f_x(1,1)=2,\qquad f_y(1,1)=2,
\]
the tangent plane is
\[
z=2+2(x-1)+2(y-1)
=
2x+2y-2.
\]

\begin{center}
\begin{tikzpicture}
\begin{axis}[
width=10cm,
height=7cm,
view={45}{25},
xlabel={$x$},
ylabel={$y$},
zlabel={$z$},
xmin=-0.4,xmax=2.3,
ymin=-0.4,ymax=2.3,
zmin=-1,zmax=8,
domain=-0.2:2.1,
y domain=-0.2:2.1,
samples=22,
samples y=22
]
\addplot3[surf,shader=interp,opacity=0.55]
{x^2+y^2};
\addplot3[surf,shader=flat,opacity=0.35,domain=0.35:1.65,y domain=0.35:1.65]
{2*x+2*y-2};
\addplot3[only marks,mark=*,mark size=2pt]
coordinates {(1,1,2)};
\end{axis}
\end{tikzpicture}
\end{center}

The plane does not replace the surface globally. Differentiability says
that, under sufficient magnification near the point of contact, the
surface and this plane become indistinguishable to first order.


\section{A Direct Differentiability Test}

Suppose we already know \(f_x(a,b)\) and \(f_y(a,b)\). To test
differentiability directly, calculate
\[
E(h,k)
=
\frac{
f(a+h,b+k)-f(a,b)
-f_x(a,b)h
-f_y(a,b)k
}{
\sqrt{h^2+k^2}
}.
\]

If
\[
E(h,k)\to0
\]
as \((h,k)\to(0,0)\), then \(f\) is differentiable.

\begin{example}
Show directly that
\[
f(x,y)=x^2+y^2
\]
is differentiable at \((a,b)\).

We have
\[
f_x(a,b)=2a,
\qquad
f_y(a,b)=2b.
\]

Now
\[
\begin{aligned}
f(a+h,b+k)-f(a,b)
&=
(a+h)^2+(b+k)^2-a^2-b^2\\
&=
2ah+2bk+h^2+k^2.
\end{aligned}
\]

Subtract the linear part:
\[
2ah+2bk.
\]

The remainder is
\[
h^2+k^2.
\]

Therefore
\[
E(h,k)
=
\frac{h^2+k^2}{\sqrt{h^2+k^2}}
=
\sqrt{h^2+k^2}.
\]

Hence
\[
E(h,k)\to0.
\]

Thus \(f\) is differentiable at every point.
\end{example}

\section{Additional Applications of Young's and Schwarz's Theorems}

\begin{example}
Let
\[
f(x,y)=x^2y^3+\sin(x+y).
\]

We have
\[
f_x=2xy^3+\cos(x+y),
\]
so
\[
f_{xy}
=
6xy^2-\sin(x+y).
\]

Also,
\[
f_y=3x^2y^2+\cos(x+y),
\]
so
\[
f_{yx}
=
6xy^2-\sin(x+y).
\]

The mixed derivatives are continuous everywhere. Therefore Schwarz's
theorem applies and
\[
f_{xy}=f_{yx}.
\]
\end{example}

\section{Additional Detailed Solved Problems}

\subsection*{Problem 3. Linear approximation}

Use linear approximation to estimate
\[
f(1.02,0.99)
\]
for
\[
f(x,y)=\sqrt{x^2+y^2}.
\]

Choose the nearby point \((1,1)\). Then
\[
f(1,1)=\sqrt2.
\]

Also,
\[
f_x=\frac{x}{\sqrt{x^2+y^2}},
\qquad
f_y=\frac{y}{\sqrt{x^2+y^2}}.
\]

Hence
\[
f_x(1,1)
=
f_y(1,1)
=
\frac1{\sqrt2}.
\]

Now
\[
h=0.02,\qquad k=-0.01.
\]

Therefore
\[
\begin{aligned}
f(1.02,0.99)
&\approx
\sqrt2
+
\frac{0.02}{\sqrt2}
-
\frac{0.01}{\sqrt2}\\
&=
\sqrt2+\frac{0.01}{\sqrt2}.
\end{aligned}
\]

\subsection*{Problem 4. Continuity does not imply differentiability}

Consider
\[
f(x,y)=\sqrt{x^2+y^2}.
\]

The function is continuous at \((0,0)\), because
\[
\sqrt{x^2+y^2}\to0.
\]

However,
\[
f_x(0,0)
=
\lim_{h\to0}
\frac{|h|}{h},
\]
which does not exist.

Hence \(f\) is not differentiable at the origin.

This is an important contrast with the theorem
\[
\text{differentiable}\Longrightarrow\text{continuous}.
\]



\section{Why Differentiability Is Stronger Than Having Partial Derivatives}

Partial derivatives test only two special directions:
\[
(h,0)
\quad\text{and}\quad
(0,k).
\]

Differentiability, however, requires one linear approximation to work
for \emph{all} small increments \((h,k)\), regardless of direction.

That is why existence of \(f_x\) and \(f_y\) is not enough.

The definition
\[
f(a+h,b+k)
=
f(a,b)
+
f_x(a,b)h
+
f_y(a,b)k
+
R(h,k)
\]
requires
\[
\frac{R(h,k)}{\sqrt{h^2+k^2}}
\to0.
\]

Thus the error must be small compared with the distance
\[
\sqrt{h^2+k^2}.
\]

\section{Little-\(o\) Notation Explained}

The notation
\[
R(h,k)=o\!\left(\sqrt{h^2+k^2}\right)
\]
means precisely that
\[
\frac{R(h,k)}
{\sqrt{h^2+k^2}}
\to0.
\]

It does not merely mean that \(R(h,k)\to0\). The remainder must go to
zero \emph{faster} than the distance from \((a+h,b+k)\) to \((a,b)\).

For example,
\[
h^2+k^2
=
o\!\left(\sqrt{h^2+k^2}\right),
\]
because
\[
\frac{h^2+k^2}{\sqrt{h^2+k^2}}
=
\sqrt{h^2+k^2}
\to0.
\]

But
\[
\sqrt{h^2+k^2}
\]
is not little-\(o\) of itself, because the ratio equals \(1\).

\section{Differentiability of a Polynomial: Detailed Verification}

Let
\[
f(x,y)=x^2+xy+y^2.
\]

We verify differentiability at an arbitrary point \((a,b)\).

First,
\[
f_x=2x+y,
\qquad
f_y=x+2y.
\]

Thus
\[
f_x(a,b)=2a+b,
\]
and
\[
f_y(a,b)=a+2b.
\]

Now
\[
\begin{aligned}
f(a+h,b+k)
={}&
(a+h)^2
+
(a+h)(b+k)
+
(b+k)^2.
\end{aligned}
\]

Expanding,
\[
\begin{aligned}
f(a+h,b+k)
={}&
a^2+2ah+h^2
+
ab+ak+bh+hk\\
&+
b^2+2bk+k^2.
\end{aligned}
\]

Subtract
\[
f(a,b)=a^2+ab+b^2.
\]

Then
\[
\begin{aligned}
\Delta f
={}&
(2a+b)h
+
(a+2b)k
+
h^2+hk+k^2.
\end{aligned}
\]

The linear part is
\[
f_x(a,b)h+f_y(a,b)k.
\]

Hence
\[
R(h,k)=h^2+hk+k^2.
\]

We estimate
\[
|hk|
\le
\frac{h^2+k^2}{2}.
\]

Therefore
\[
|R(h,k)|
\le
h^2+\frac{h^2+k^2}{2}+k^2
=
\frac32(h^2+k^2).
\]

Thus
\[
\frac{|R(h,k)|}{\sqrt{h^2+k^2}}
\le
\frac32\sqrt{h^2+k^2}.
\]

The right-hand side tends to \(0\). Hence
\[
f
\]
is differentiable at \((a,b)\).

Since \((a,b)\) was arbitrary, \(f\) is differentiable everywhere.

\section{Deriving the Tangent Plane Formula from Differentiability}

Differentiability gives
\[
f(a+h,b+k)
\approx
f(a,b)
+
f_x(a,b)h
+
f_y(a,b)k.
\]

Since
\[
h=x-a,
\qquad
k=y-b,
\]
we obtain
\[
z
\approx
f(a,b)
+
f_x(a,b)(x-a)
+
f_y(a,b)(y-b).
\]

The right-hand side is a plane. Therefore the tangent plane is not an
independent formula to memorize; it is exactly the linear part of the
definition of differentiability.

\section{Total Differential}

The linear change
\[
df
=
f_x\,dx+f_y\,dy
\]
is called the total differential.

For small changes,
\[
\Delta f
\approx df.
\]

If
\[
z=f(x,y),
\]
then
\[
dz=f_x\,dx+f_y\,dy.
\]

\begin{example}
Let
\[
z=x^2y+y^2.
\]

Then
\[
z_x=2xy,
\qquad
z_y=x^2+2y.
\]

Therefore
\[
dz
=
2xy\,dx
+
(x^2+2y)\,dy.
\]
\end{example}

\section{Error in Linear Approximation}

Suppose
\[
L(x,y)
=
f(a,b)
+
f_x(a,b)(x-a)
+
f_y(a,b)(y-b).
\]

Then the approximation error is
\[
E(x,y)=f(x,y)-L(x,y).
\]

Differentiability means
\[
\frac{|E(x,y)|}
{\sqrt{(x-a)^2+(y-b)^2}}
\to0.
\]

Thus the error becomes negligible compared with the distance from the
base point.

\section{Another Example: Partial Derivatives Exist but Differentiability Fails}

Define
\[
f(x,y)
=
\begin{cases}
\dfrac{x^2y}{x^4+y^2},&(x,y)\ne(0,0),\\
0,&(x,y)=(0,0).
\end{cases}
\]

Along both coordinate axes,
\[
f(h,0)=0,
\qquad
f(0,k)=0.
\]

Hence
\[
f_x(0,0)=0,
\qquad
f_y(0,0)=0.
\]

However, along
\[
y=x^2,
\]
\[
f(x,x^2)=\frac12.
\]

The function is not even continuous at the origin. Therefore it cannot
be differentiable there.

This example reinforces the hierarchy:
\[
\text{partial derivatives}
\]
look only in coordinate directions, while
\[
\text{differentiability}
\]
controls all directions simultaneously.

\section{A Practical Differentiability Checklist}

At a point \((a,b)\):

\begin{enumerate}
\item Check whether the function is continuous. If not, it is not
differentiable.

\item Compute \(f_x(a,b)\) and \(f_y(a,b)\).

\item If \(f_x,f_y\) exist in a neighbourhood and are continuous at
\((a,b)\), differentiability follows from the sufficient theorem.

\item If continuity of the partial derivatives is unavailable, use the
definition directly.

\item In the direct test, subtract the linear part and divide the
remainder by
\[
\sqrt{h^2+k^2}.
\]
\end{enumerate}

\section{Young's and Schwarz's Theorems: Interpretation}

The equality
\[
f_{xy}=f_{yx}
\]
says that, under suitable smoothness conditions, the order in which we
differentiate with respect to \(x\) and \(y\) does not matter.

This is analogous to the fact that sufficiently smooth processes can
be performed in either order without changing the final result.

But the hypotheses are essential. Without suitable continuity or
regularity, mixed partial derivatives may exist and still be unequal.

\section{More Detailed Solved Problems}

\subsection*{Problem 5. Test differentiability directly}

Define
\[
f(x,y)
=
\begin{cases}
\dfrac{x^2y}{\sqrt{x^2+y^2}},&(x,y)\ne(0,0),\\
0,&(x,y)=(0,0).
\end{cases}
\]

First,
\[
f(h,0)=0,\qquad f(0,k)=0,
\]
so
\[
f_x(0,0)=f_y(0,0)=0.
\]

For differentiability we must test
\[
\frac{|f(h,k)|}{\sqrt{h^2+k^2}}.
\]

Now
\[
\frac{|f(h,k)|}{\sqrt{h^2+k^2}}
=
\frac{|h|^2|k|}{h^2+k^2}.
\]

Let
\[
r=\sqrt{h^2+k^2}.
\]

Since
\[
|h|\le r,\qquad |k|\le r,
\]
we get
\[
\frac{|h|^2|k|}{h^2+k^2}
\le
\frac{r^3}{r^2}
=
r.
\]

Thus the expression tends to zero. Therefore \(f\) is differentiable
at the origin.

\subsection*{Problem 6. Tangent plane and approximation}

Let
\[
f(x,y)=e^{x+y}.
\]

At \((0,0)\),
\[
f(0,0)=1,
\]
\[
f_x=e^{x+y},
\qquad
f_y=e^{x+y}.
\]

Hence
\[
f_x(0,0)=f_y(0,0)=1.
\]

The tangent plane is
\[
z=1+x+y.
\]

Therefore, near the origin,
\[
e^{x+y}\approx1+x+y.
\]

For example,
\[
e^{0.02-0.01}
=
e^{0.01}
\approx
1.01.
\]


\section{Detailed Solved Problems}

\subsection*{Problem 1. Prove differentiability}
Show that
\[
f(x,y)=x^2+xy+y^2
\]
is differentiable everywhere.

Its first partial derivatives are
\[
f_x=2x+y,
\qquad
f_y=x+2y.
\]
These are polynomials, hence continuous everywhere. Therefore by the
standard sufficient condition, \(f\) is differentiable at every
point of \(\R^2\).

\subsection*{Problem 2. A harmonic-function application}
Let
\[
u(x,y)=e^x\cos y.
\]
Then
\[
u_x=e^x\cos y,
\qquad
u_{xx}=e^x\cos y.
\]
Also
\[
u_y=-e^x\sin y,
\qquad
u_{yy}=-e^x\cos y.
\]
Therefore
\[
u_{xx}+u_{yy}=0.
\]

\section{Graded Practice Problems}
\begin{enumerate}
\item Prove directly that differentiability implies continuity.
\item Determine whether
\[
f(x,y)=
\begin{cases}
\dfrac{x^2y}{x^2+y^2},&(x,y)\ne(0,0),\\
0,&(0,0)
\end{cases}
\]
is differentiable at the origin.
\item Verify Schwarz's theorem for
\[
f(x,y)=x^2y^3+\sin(xy).
\]
\item Give an example showing that mixed partial derivatives can
differ when the continuity hypotheses fail.
\end{enumerate}

% ============================================================
\chapter{Derivatives of Composite Functions}
% ============================================================

\section*{Learning Objectives}
Students should be able to:
\begin{itemize}
\item identify direct and indirect dependence among variables;
\item use the multivariable chain rule;
\item differentiate when all variables depend on one parameter;
\item use dependency diagrams to avoid missing terms.
\end{itemize}

\section*{Historical Motivation}

Many physical quantities depend on intermediate variables.
Temperature may depend on position, while position depends on time.
The chain rule organizes these layers of dependence and tells us how
small changes propagate through a system.

\section{Two Intermediate Variables}

Suppose
\[
z=f(u,v),
\]
where
\[
u=u(x,y),
\qquad
v=v(x,y).
\]

\begin{theorem}[Chain rule]
If the functions involved are differentiable, then
\[
\pd{z}{x}
=
\pd{z}{u}\pd{u}{x}
+
\pd{z}{v}\pd{v}{x},
\]
and
\[
\pd{z}{y}
=
\pd{z}{u}\pd{u}{y}
+
\pd{z}{v}\pd{v}{y}.
\]
\end{theorem}

\section{Derivation from Differentials}

For differentiable \(z=f(u,v)\),
\[
dz=f_u\,du+f_v\,dv.
\]
Now
\[
du=u_x\,dx+u_y\,dy,
\]
\[
dv=v_x\,dx+v_y\,dy.
\]
Substitution gives
\[
\begin{aligned}
dz
&=
f_u(u_x\,dx+u_y\,dy)
+
f_v(v_x\,dx+v_y\,dy)\\
&=
(f_u u_x+f_v v_x)\,dx
+
(f_u u_y+f_v v_y)\,dy.
\end{aligned}
\]
But
\[
dz=z_x\,dx+z_y\,dy.
\]
Comparing coefficients,
\[
z_x=f_u u_x+f_v v_x,
\]
\[
z_y=f_u u_y+f_v v_y.
\]

\section{Dependency Diagram}

\begin{center}
\begin{tikzpicture}[
node distance=1.5cm and 1.9cm,
every node/.style={draw,rounded corners,minimum width=1cm},
>=Stealth
]
\node (x) at (0,1) {$x$};
\node (y) at (0,-1) {$y$};
\node (u) at (3,1) {$u$};
\node (v) at (3,-1) {$v$};
\node (z) at (6,0) {$z$};

\draw[->] (x)--(u);
\draw[->] (x)--(v);
\draw[->] (y)--(u);
\draw[->] (y)--(v);
\draw[->] (u)--(z);
\draw[->] (v)--(z);
\end{tikzpicture}
\end{center}

For a derivative with respect to \(x\), add the products of
derivatives along every path from \(x\) to \(z\).

\section{Detailed Example}

Let
\[
z=u^2+v^2,
\qquad
u=x+y,
\qquad
v=x-y.
\]
First,
\[
z_u=2u,\qquad z_v=2v.
\]
Since
\[
u_x=1,\qquad v_x=1,
\]
we have
\[
z_x=2u+2v.
\]
Substituting \(u=x+y,v=x-y\),
\[
z_x=4x.
\]

Likewise
\[
u_y=1,\qquad v_y=-1,
\]
so
\[
z_y=2u-2v=4y.
\]

\section{One-Parameter Chain Rule}

If
\[
z=f(x,y),
\qquad
x=x(t),\quad y=y(t),
\]
then
\[
\frac{dz}{dt}
=
f_x\frac{dx}{dt}
+
f_y\frac{dy}{dt}.
\]

\begin{example}
Let
\[
z=x^2y,
\qquad
x=t^2,
\qquad
y=\sin t.
\]
Then
\[
z_x=2xy,\qquad z_y=x^2,
\]
and
\[
\frac{dx}{dt}=2t,
\qquad
\frac{dy}{dt}=\cos t.
\]
Therefore
\[
\begin{aligned}
\frac{dz}{dt}
&=
(2xy)(2t)+x^2\cos t\\
&=
4t^3\sin t+t^4\cos t.
\end{aligned}
\]

As a check, direct substitution gives
\[
z=t^4\sin t,
\]
whose derivative is the same expression.
\end{example}


\section{Derivation of the Chain Rule from Differentiability}

Suppose
\[
z=f(u,v),
\]
where
\[
u=u(x,y),
\qquad
v=v(x,y).
\]

Let \(x\) change by \(\Delta x\), while \(y\) is held fixed. Then
\[
u\to u+\Delta u,
\qquad
v\to v+\Delta v.
\]

Since \(f\) is differentiable,
\[
\Delta z
=
f_u\Delta u
+
f_v\Delta v
+
\text{higher-order error}.
\]

Also, since \(u\) and \(v\) are differentiable,
\[
\Delta u
=
u_x\Delta x
+
o(\Delta x),
\]
and
\[
\Delta v
=
v_x\Delta x
+
o(\Delta x).
\]

Substituting,
\[
\Delta z
=
f_u u_x\Delta x
+
f_v v_x\Delta x
+
o(\Delta x).
\]

Divide by \(\Delta x\):
\[
\frac{\Delta z}{\Delta x}
=
f_u u_x
+
f_v v_x
+
o(1).
\]

Letting \(\Delta x\to0\),
\[
z_x
=
f_u u_x
+
f_v v_x.
\]

Similarly,
\[
z_y
=
f_u u_y
+
f_v v_y.
\]

This derivation explains why each path of dependence contributes one
product of derivatives.

\section{Chain Rule with One Intermediate Variable}

If
\[
z=f(u),
\qquad
u=u(x,y),
\]
then
\[
z_x=f'(u)u_x,
\qquad
z_y=f'(u)u_y.
\]

\begin{example}
Let
\[
z=\log(x^2+y^2).
\]

Set
\[
u=x^2+y^2.
\]
Then
\[
z=\log u.
\]

Hence
\[
z_x
=
\frac1u(2x)
=
\frac{2x}{x^2+y^2},
\]
and
\[
z_y
=
\frac{2y}{x^2+y^2}.
\]
\end{example}

\section{Several Intermediate Variables}

The same principle extends immediately. If
\[
z=f(u_1,u_2,\ldots,u_m)
\]
and each \(u_i\) depends on \(x\), then
\[
\frac{\partial z}{\partial x}
=
\sum_{i=1}^m
\frac{\partial z}{\partial u_i}
\frac{\partial u_i}{\partial x}.
\]

The formula simply adds the contributions from all dependency paths.

\section{Additional Detailed Solved Problems}

\subsection*{Problem 3. Two intermediate variables}

Let
\[
z=u^2v+e^v,
\qquad
u=x^2+y,
\qquad
v=x-y^2.
\]

First,
\[
z_u=2uv,
\qquad
z_v=u^2+e^v.
\]

Also,
\[
u_x=2x,\qquad v_x=1.
\]

Therefore
\[
\begin{aligned}
z_x
&=
z_u u_x+z_v v_x\\
&=
(2uv)(2x)+(u^2+e^v)\\
&=
4xuv+u^2+e^v.
\end{aligned}
\]

For \(y\),
\[
u_y=1,\qquad v_y=-2y.
\]

Hence
\[
\begin{aligned}
z_y
&=
z_u u_y+z_v v_y\\
&=
2uv-2y(u^2+e^v).
\end{aligned}
\]

Finally substitute
\[
u=x^2+y,\qquad v=x-y^2
\]
if an answer entirely in \(x,y\) is required.

\subsection*{Problem 4. Parameterized motion}

Let
\[
z=x^2+xy+y^2,
\qquad
x=e^t,
\qquad
y=t^2.
\]

Then
\[
z_x=2x+y,
\qquad
z_y=x+2y.
\]

Also,
\[
\frac{dx}{dt}=e^t,
\qquad
\frac{dy}{dt}=2t.
\]

Therefore
\[
\frac{dz}{dt}
=
(2x+y)e^t
+
(x+2y)2t.
\]

Substituting
\[
x=e^t,\qquad y=t^2,
\]
gives the derivative entirely as a function of \(t\).



\section{Why the Multivariable Chain Rule Has a Sum}

Suppose
\[
z=f(u,v),
\]
and both \(u\) and \(v\) depend on \(x\).

A change in \(x\) affects \(z\) in two ways:
\[
x\longrightarrow u\longrightarrow z,
\]
and
\[
x\longrightarrow v\longrightarrow z.
\]

The first contribution is
\[
f_u u_x,
\]
and the second is
\[
f_v v_x.
\]

Since both effects occur simultaneously, they add:
\[
z_x=f_u u_x+f_v v_x.
\]

This is the conceptual reason for the sum in the chain rule.

\section{Tree Method for the Chain Rule}

A dependency tree can be used to avoid missing terms.

If
\[
z=f(u,v),
\qquad
u=u(x,y),
\qquad
v=v(x,y),
\]
then for \(z_x\) list all paths from \(x\) to \(z\):

\[
x\to u\to z,
\qquad
x\to v\to z.
\]

Multiply derivatives along each path and add:
\[
z_x
=
u_xz_u+v_xz_v.
\]

For \(z_y\):
\[
y\to u\to z,
\qquad
y\to v\to z,
\]
so
\[
z_y
=
u_yz_u+v_yz_v.
\]

\section{Chain Rule with a Parameter: Detailed Derivation}

Let
\[
z=f(x,y),
\]
where
\[
x=x(t),\qquad y=y(t).
\]

For a small change \(\Delta t\),
\[
\Delta x\approx x'(t)\Delta t,
\qquad
\Delta y\approx y'(t)\Delta t.
\]

Differentiability of \(f\) gives
\[
\Delta z
\approx
f_x\Delta x+f_y\Delta y.
\]

Therefore
\[
\Delta z
\approx
\left(
f_xx'(t)+f_yy'(t)
\right)\Delta t.
\]

Dividing by \(\Delta t\) and passing to the limit,
\[
\frac{dz}{dt}
=
f_x\frac{dx}{dt}
+
f_y\frac{dy}{dt}.
\]

\section{Second Derivatives of Composite Functions}

Suppose
\[
z=f(u),
\qquad
u=u(x).
\]

Then
\[
z_x=f'(u)u_x.
\]

Differentiate again:
\[
z_{xx}
=
f''(u)(u_x)^2
+
f'(u)u_{xx}.
\]

In several variables, repeated application of the chain rule produces
similar terms.

\begin{example}
Let
\[
z=e^u,
\qquad
u=x^2+y^2.
\]

Then
\[
z_x=e^u(2x).
\]

Differentiate again:
\[
z_{xx}
=
2e^u
+
4x^2e^u.
\]

Since \(u=x^2+y^2\),
\[
z_{xx}
=
(2+4x^2)e^{x^2+y^2}.
\]
\end{example}

\section{More Detailed Solved Problems}

\subsection*{Problem 5. Nested composite function}

Let
\[
z=\sin(u^2+v),
\qquad
u=x+y,
\qquad
v=xy.
\]

Set
\[
w=u^2+v.
\]

Then
\[
z=\sin w,
\]
so
\[
z_w=\cos w.
\]

Also,
\[
w_u=2u,
\qquad
w_v=1.
\]

Now
\[
u_x=1,\qquad v_x=y.
\]

Thus
\[
\begin{aligned}
z_x
&=
z_w(w_u u_x+w_v v_x)\\
&=
\cos(u^2+v)(2u+y).
\end{aligned}
\]

Similarly,
\[
u_y=1,\qquad v_y=x,
\]
so
\[
z_y
=
\cos(u^2+v)(2u+x).
\]

Finally substitute
\[
u=x+y,\qquad v=xy.
\]

\subsection*{Problem 6. Motion along a curve}

Suppose
\[
T(x,y)=x^2+y^2
\]
is a temperature field and a particle moves according to
\[
x=t,\qquad y=t^2.
\]

Then
\[
T_x=2x,\qquad T_y=2y.
\]

Also,
\[
\frac{dx}{dt}=1,
\qquad
\frac{dy}{dt}=2t.
\]

Therefore
\[
\frac{dT}{dt}
=
2x+4ty.
\]

Substituting
\[
x=t,\qquad y=t^2,
\]
gives
\[
\frac{dT}{dt}
=
2t+4t^3.
\]


\section{Detailed Solved Problems}

\subsection*{Problem 1}
Let
\[
z=e^{uv},
\qquad
u=x+y,
\qquad
v=x-y.
\]
Then
\[
z_u=ve^{uv},
\qquad
z_v=ue^{uv}.
\]
Thus
\[
z_x
=
ve^{uv}+ue^{uv}
=
(u+v)e^{uv}
=
2xe^{uv}.
\]
Similarly,
\[
z_y
=
ve^{uv}-ue^{uv}
=
(v-u)e^{uv}
=
-2ye^{uv}.
\]

\subsection*{Problem 2}
Let
\[
z=\log(u^2+v^2),
\quad
u=x^2-y^2,
\quad
v=2xy.
\]
Then
\[
z_u=\frac{2u}{u^2+v^2},
\qquad
z_v=\frac{2v}{u^2+v^2}.
\]
Also
\[
u_x=2x,\quad v_x=2y,
\qquad
u_y=-2y,\quad v_y=2x.
\]
Hence
\[
z_x
=
\frac{4xu+4yv}{u^2+v^2},
\]
and
\[
z_y
=
\frac{-4yu+4xv}{u^2+v^2}.
\]

\section{Graded Practice Problems}
\begin{enumerate}
\item If
\[
z=u^3+v^3,\quad u=x+y,\quad v=x-y,
\]
find \(z_x,z_y\).
\item If
\[
z=e^{xy},\quad x=t^2,\quad y=\cos t,
\]
find \(dz/dt\).
\item Draw a dependency diagram and derive the chain rule when
\(z=f(u,v,w)\) and each of \(u,v,w\) depends on \(x,y\).
\end{enumerate}

% ============================================================
\chapter{Change of Variables and Jacobians}
% ============================================================

\section*{Learning Objectives}
Students should be able to:
\begin{itemize}
\item understand why variables are changed;
\item calculate a Jacobian determinant;
\item derive transformation formulae for first derivatives;
\item handle a linear transformation;
\item transform derivatives between Cartesian and polar coordinates.
\end{itemize}

\section*{Historical Motivation}

Changing coordinates can reveal the geometry hidden in an algebraic
problem. Circular symmetry suggests polar coordinates, while other
geometries suggest other transformations. Jacobi's determinant gives a
systematic way to describe how variables change together.

\section{Change of Independent Variables}

Suppose
\[
u=u(x,y),
\qquad
v=v(x,y).
\]
If the transformation is invertible, we may also write
\[
x=x(u,v),
\qquad
y=y(u,v).
\]

A function
\[
z=f(x,y)
\]
may then be regarded as
\[
z=F(u,v).
\]


\section{Visualizing a Change of Coordinates}

For the linear transformation
\[
u=x+y,\qquad v=x-y,
\]
lines of constant \(u\) and constant \(v\) form a rotated coordinate
grid in the original \(xy\)-plane.

\begin{center}
\begin{tikzpicture}[scale=0.9]
% xy grid
\begin{scope}[xshift=-4.2cm]
\draw[->] (-2.7,0)--(2.7,0) node[right] {$x$};
\draw[->] (0,-2.7)--(0,2.7) node[above] {$y$};
\foreach \c in {-2,-1,1,2}{
  \draw[dashed] (-2.4,{\c+2.4})--(2.4,{\c-2.4});
  \draw[dotted] (-2.4,{-\c-2.4})--(2.4,{-\c+2.4});
}
\node at (0,-3.2) {constant \(u\) and \(v\) lines in the \(xy\)-plane};
\end{scope}

\draw[->,very thick] (-0.8,0)--(0.8,0);

% uv grid
\begin{scope}[xshift=4.2cm]
\draw[->] (-2.7,0)--(2.7,0) node[right] {$u$};
\draw[->] (0,-2.7)--(0,2.7) node[above] {$v$};
\foreach \c in {-2,-1,1,2}{
  \draw[dashed] (\c,-2.4)--(\c,2.4);
  \draw[dotted] (-2.4,\c)--(2.4,\c);
}
\node at (0,-3.2) {the same families become coordinate lines};
\end{scope}
\end{tikzpicture}
\end{center}

A change of variables is therefore not merely a symbolic
substitution: it replaces one coordinate grid by another one chosen to
fit the geometry of the problem more naturally.


\section{The Jacobian}

\begin{definition}
The Jacobian of \(u,v\) with respect to \(x,y\) is
\[
\frac{\partial(u,v)}{\partial(x,y)}
=
\begin{vmatrix}
u_x&u_y\\
v_x&v_y
\end{vmatrix}
=
u_xv_y-u_yv_x.
\]
\end{definition}

A nonzero Jacobian is the elementary signal that the transformation is
locally invertible.

\section{A Linear Transformation}

Let
\[
u=x+y,
\qquad
v=x-y.
\]
Then
\[
u_x=1,\quad u_y=1,
\qquad
v_x=1,\quad v_y=-1.
\]
Therefore
\[
\frac{\partial(u,v)}{\partial(x,y)}
=
\begin{vmatrix}
1&1\\
1&-1
\end{vmatrix}
=-2.
\]

Solving the equations,
\[
x=\frac{u+v}{2},
\qquad
y=\frac{u-v}{2}.
\]

\section{Transformation of Differential Operators}

If \(z=F(u,v)\), then
\[
z_x=F_u u_x+F_v v_x.
\]
Symbolically,
\[
\frac{\partial}{\partial x}
=
u_x\frac{\partial}{\partial u}
+
v_x\frac{\partial}{\partial v}.
\]
Similarly,
\[
\frac{\partial}{\partial y}
=
u_y\frac{\partial}{\partial u}
+
v_y\frac{\partial}{\partial v}.
\]

For
\[
u=x+y,\qquad v=x-y,
\]
we obtain
\[
\frac{\partial}{\partial x}
=
\frac{\partial}{\partial u}
+
\frac{\partial}{\partial v},
\]
and
\[
\frac{\partial}{\partial y}
=
\frac{\partial}{\partial u}
-
\frac{\partial}{\partial v}.
\]

\section{A Second-Order Application}

Using the preceding formulas,
\[
z_{xx}
=
z_{uu}+2z_{uv}+z_{vv},
\]
and
\[
z_{yy}
=
z_{uu}-2z_{uv}+z_{vv}.
\]
Subtracting,
\[
z_{xx}-z_{yy}=4z_{uv}.
\]

\section{Reciprocal Jacobian}

If the transformation is invertible and
\[
J=
\frac{\partial(u,v)}{\partial(x,y)}
\ne0,
\]
then
\[
\frac{\partial(x,y)}{\partial(u,v)}
=
\frac1J.
\]

For the linear transformation above,
\[
\frac{\partial(x,y)}{\partial(u,v)}
=
-\frac12,
\]
which is indeed the reciprocal of \(-2\).

\section{Polar Coordinates}

Let
\[
x=r\cos\theta,
\qquad
y=r\sin\theta.
\]
If
\[
F(r,\theta)=f(x,y),
\]
then by the chain rule
\[
F_r
=
f_x\cos\theta+f_y\sin\theta.
\]

Also,
\[
x_\theta=-r\sin\theta,
\qquad
y_\theta=r\cos\theta,
\]
so
\[
F_\theta
=
-rf_x\sin\theta
+
rf_y\cos\theta.
\]

Solving these two equations for \(f_x,f_y\),
\[
f_x
=
F_r\cos\theta
-
\frac{\sin\theta}{r}F_\theta,
\]
\[
f_y
=
F_r\sin\theta
+
\frac{\cos\theta}{r}F_\theta.
\]

\section{Jacobian of Polar Coordinates}

For
\[
x=r\cos\theta,\qquad y=r\sin\theta,
\]
\[
x_r=\cos\theta,
\quad
x_\theta=-r\sin\theta,
\]
\[
y_r=\sin\theta,
\quad
y_\theta=r\cos\theta.
\]
Thus
\[
\begin{aligned}
\frac{\partial(x,y)}{\partial(r,\theta)}
&=
\begin{vmatrix}
\cos\theta&-r\sin\theta\\
\sin\theta&r\cos\theta
\end{vmatrix}\\
&=
r\cos^2\theta+r\sin^2\theta\\
&=r.
\end{aligned}
\]


\section{General Change-of-Variable Formulae}

Suppose
\[
u=u(x,y),
\qquad
v=v(x,y),
\]
and
\[
z=F(u,v).
\]

Then
\[
z_x=F_u u_x+F_v v_x,
\]
and
\[
z_y=F_u u_y+F_v v_y.
\]

In matrix form,
\[
\begin{pmatrix}
z_x\\
z_y
\end{pmatrix}
=
\begin{pmatrix}
u_x&v_x\\
u_y&v_y
\end{pmatrix}
\begin{pmatrix}
F_u\\
F_v
\end{pmatrix}.
\]

The determinant
\[
J
=
\frac{\partial(u,v)}{\partial(x,y)}
=
u_xv_y-u_yv_x
\]
measures whether the transformation is locally invertible.

If
\[
J\ne0,
\]
we can solve locally for \(x,y\) in terms of \(u,v\).

\section{Jacobians and Local Area Scaling}

Besides its algebraic role, the absolute value
\[
|J|
\]
measures the local area-scaling factor of the transformation.

For example, for
\[
u=2x,\qquad v=3y,
\]
we have
\[
J
=
\begin{vmatrix}
2&0\\
0&3
\end{vmatrix}
=6.
\]

A small rectangle of area \(\Delta x\,\Delta y\) is therefore mapped
approximately to a rectangle of area
\[
6\,\Delta x\,\Delta y.
\]

This geometric interpretation becomes especially important later in
multiple integration.

\section{A Useful Linear Change of Variables}

Let
\[
u=x+y,\qquad v=x-y.
\]

Then
\[
x=\frac{u+v}{2},
\qquad
y=\frac{u-v}{2}.
\]

The derivative operators transform as
\[
\partial_x=\partial_u+\partial_v,
\]
and
\[
\partial_y=\partial_u-\partial_v.
\]

Hence
\[
\partial_{xx}
=
\partial_{uu}
+
2\partial_{uv}
+
\partial_{vv},
\]
while
\[
\partial_{yy}
=
\partial_{uu}
-
2\partial_{uv}
+
\partial_{vv}.
\]

Therefore
\[
\partial_{xx}-\partial_{yy}
=
4\partial_{uv}.
\]

This is a standard example of how an appropriate change of variables
can simplify a second-order differential expression.

\section{More Detail on Polar Coordinates}

With
\[
x=r\cos\theta,
\qquad
y=r\sin\theta,
\]
we have
\[
r=\sqrt{x^2+y^2},
\qquad
\theta=\tan^{-1}\frac{y}{x}
\]
locally.

The Jacobian is
\[
\frac{\partial(x,y)}{\partial(r,\theta)}
=
\begin{vmatrix}
\cos\theta&-r\sin\theta\\
\sin\theta&r\cos\theta
\end{vmatrix}
=r.
\]

Thus, away from \(r=0\),
\[
\frac{\partial(r,\theta)}{\partial(x,y)}
=
\frac1r.
\]

From
\[
F(r,\theta)
=
f(r\cos\theta,r\sin\theta),
\]
we obtain
\[
F_r=f_x\cos\theta+f_y\sin\theta,
\]
and
\[
F_\theta
=
-rf_x\sin\theta
+
rf_y\cos\theta.
\]

Solving,
\[
f_x
=
F_r\cos\theta
-
\frac{\sin\theta}{r}F_\theta,
\]
and
\[
f_y
=
F_r\sin\theta
+
\frac{\cos\theta}{r}F_\theta.
\]

\section{Additional Detailed Solved Problems}

\subsection*{Problem 3. Reciprocal Jacobian}

Let
\[
u=x+y,\qquad v=x-y.
\]

We already found
\[
\frac{\partial(u,v)}{\partial(x,y)}
=-2.
\]

From
\[
x=\frac{u+v}{2},
\qquad
y=\frac{u-v}{2},
\]
we have
\[
x_u=\frac12,\quad x_v=\frac12,
\]
and
\[
y_u=\frac12,\quad y_v=-\frac12.
\]

Therefore
\[
\frac{\partial(x,y)}{\partial(u,v)}
=
\begin{vmatrix}
1/2&1/2\\
1/2&-1/2
\end{vmatrix}
=
-\frac12.
\]

Thus
\[
\frac{\partial(u,v)}{\partial(x,y)}
\frac{\partial(x,y)}{\partial(u,v)}
=
(-2)\left(-\frac12\right)
=1.
\]

\subsection*{Problem 4. A nonlinear Jacobian}

Let
\[
u=x^2-y^2,
\qquad
v=2xy.
\]

Then
\[
u_x=2x,\quad u_y=-2y,
\]
and
\[
v_x=2y,\quad v_y=2x.
\]

Therefore
\[
\begin{aligned}
\frac{\partial(u,v)}{\partial(x,y)}
&=
\begin{vmatrix}
2x&-2y\\
2y&2x
\end{vmatrix}\\
&=
4x^2+4y^2\\
&=
4(x^2+y^2).
\end{aligned}
\]

The transformation fails to be locally invertible only at the origin,
where the Jacobian vanishes.



\section{Why We Change Variables}

A change of variables is useful when the original coordinates hide the
structure of a problem.

For example, the expression
\[
x^2+y^2
\]
is naturally adapted to polar coordinates because
\[
x^2+y^2=r^2.
\]

Similarly, the expressions
\[
x+y
\quad\text{and}\quad
x-y
\]
suggest the transformation
\[
u=x+y,\qquad v=x-y.
\]

A good change of variables simplifies either the formula, the geometry,
or the differential operators.

\section{Jacobian as the Determinant of the Differential Transformation}

For
\[
u=u(x,y),
\qquad
v=v(x,y),
\]
small changes satisfy approximately
\[
\Delta u
\approx
u_x\Delta x+u_y\Delta y,
\]
\[
\Delta v
\approx
v_x\Delta x+v_y\Delta y.
\]

In matrix form,
\[
\begin{pmatrix}
\Delta u\\
\Delta v
\end{pmatrix}
\approx
\begin{pmatrix}
u_x&u_y\\
v_x&v_y
\end{pmatrix}
\begin{pmatrix}
\Delta x\\
\Delta y
\end{pmatrix}.
\]

The determinant of this matrix is
\[
J
=
u_xv_y-u_yv_x.
\]

Thus the Jacobian measures whether the local transformation collapses
area.

If
\[
J=0,
\]
the linearized transformation is singular.

If
\[
J\ne0,
\]
the transformation is locally invertible under suitable regularity
conditions.

\section{Reciprocal Jacobian in Detail}

Suppose
\[
u=u(x,y),
\qquad
v=v(x,y),
\]
and the inverse transformation
\[
x=x(u,v),
\qquad
y=y(u,v)
\]
exists.

Then
\[
\frac{\partial(u,v)}{\partial(x,y)}
\frac{\partial(x,y)}{\partial(u,v)}
=1.
\]

This is the determinant version of the matrix identity
\[
A^{-1}A=I.
\]

\begin{example}
Let
\[
u=2x+y,
\qquad
v=x-y.
\]

Then
\[
\frac{\partial(u,v)}{\partial(x,y)}
=
\begin{vmatrix}
2&1\\
1&-1
\end{vmatrix}
=-3.
\]

Solving,
\[
x=\frac{u+v}{3},
\qquad
y=\frac{u-2v}{3}.
\]

Hence
\[
\frac{\partial(x,y)}{\partial(u,v)}
=
\begin{vmatrix}
1/3&1/3\\
1/3&-2/3
\end{vmatrix}
=
-\frac13.
\]

Therefore
\[
(-3)\left(-\frac13\right)=1.
\]
\end{example}

\section{Operator Transformation: Step-by-Step}

Let
\[
u=u(x,y),
\qquad
v=v(x,y),
\]
and let
\[
z=F(u,v).
\]

Then
\[
z_x
=
F_u u_x+F_v v_x.
\]

This can be written symbolically as
\[
\frac{\partial}{\partial x}
=
u_x\frac{\partial}{\partial u}
+
v_x\frac{\partial}{\partial v}.
\]

Similarly,
\[
\frac{\partial}{\partial y}
=
u_y\frac{\partial}{\partial u}
+
v_y\frac{\partial}{\partial v}.
\]

These operator identities should be applied to a function on the
right. They are shorthand for the chain rule.

\section{Detailed Polar-Coordinate Derivation}

Let
\[
x=r\cos\theta,\qquad y=r\sin\theta.
\]

If
\[
F(r,\theta)=f(x,y),
\]
then
\[
F_r
=
f_xx_r+f_yy_r.
\]

Now
\[
x_r=\cos\theta,
\qquad
y_r=\sin\theta,
\]
so
\[
F_r
=
f_x\cos\theta+f_y\sin\theta.
\]

Likewise,
\[
F_\theta
=
f_xx_\theta+f_yy_\theta.
\]

Since
\[
x_\theta=-r\sin\theta,
\qquad
y_\theta=r\cos\theta,
\]
we obtain
\[
F_\theta
=
-rf_x\sin\theta
+
rf_y\cos\theta.
\]

To solve for \(f_x,f_y\), divide the second equation by \(r\):
\[
\frac1rF_\theta
=
-f_x\sin\theta
+
f_y\cos\theta.
\]

Thus
\[
\begin{pmatrix}
F_r\\
\frac1rF_\theta
\end{pmatrix}
=
\begin{pmatrix}
\cos\theta&\sin\theta\\
-\sin\theta&\cos\theta
\end{pmatrix}
\begin{pmatrix}
f_x\\
f_y
\end{pmatrix}.
\]

The matrix is a rotation matrix. Its inverse is its transpose.
Therefore
\[
f_x
=
F_r\cos\theta
-
\frac{\sin\theta}{r}F_\theta,
\]
\[
f_y
=
F_r\sin\theta
+
\frac{\cos\theta}{r}F_\theta.
\]

\section{More Detailed Solved Problems}

\subsection*{Problem 5. Transform a first-order expression}

Let
\[
u=x+y,\qquad v=x-y.
\]

Show that
\[
z_x+z_y=2z_u.
\]

Since
\[
\partial_x=\partial_u+\partial_v
\]
and
\[
\partial_y=\partial_u-\partial_v,
\]
we have
\[
\begin{aligned}
z_x+z_y
&=
(z_u+z_v)+(z_u-z_v)\\
&=
2z_u.
\end{aligned}
\]

Similarly,
\[
z_x-z_y=2z_v.
\]

\subsection*{Problem 6. Polar Jacobian}

For
\[
x=r\cos\theta,\qquad y=r\sin\theta,
\]
\[
\frac{\partial(x,y)}{\partial(r,\theta)}
=
\begin{vmatrix}
\cos\theta&-r\sin\theta\\
\sin\theta&r\cos\theta
\end{vmatrix}.
\]

Hence
\[
\begin{aligned}
J
&=
r\cos^2\theta+r\sin^2\theta\\
&=
r.
\end{aligned}
\]

Thus polar coordinates become singular at
\[
r=0,
\]
which reflects the fact that every angle \(\theta\) represents the same
point at the origin.


\section{Detailed Solved Problems}

\subsection*{Problem 1. A quadratic transformation}
For
\[
u=x^2-y^2,
\qquad
v=2xy,
\]
find the Jacobian.

We have
\[
u_x=2x,\quad u_y=-2y,
\]
\[
v_x=2y,\quad v_y=2x.
\]
Therefore
\[
\begin{aligned}
\frac{\partial(u,v)}{\partial(x,y)}
&=
\begin{vmatrix}
2x&-2y\\
2y&2x
\end{vmatrix}\\
&=
4x^2+4y^2\\
&=
4(x^2+y^2).
\end{aligned}
\]

\subsection*{Problem 2. Transform a second-order expression}
Let
\[
u=x+y,\qquad v=x-y.
\]
Then
\[
\partial_x=\partial_u+\partial_v,
\qquad
\partial_y=\partial_u-\partial_v.
\]
Hence
\[
z_{xx}
=
z_{uu}+2z_{uv}+z_{vv},
\]
\[
z_{yy}
=
z_{uu}-2z_{uv}+z_{vv}.
\]
Therefore
\[
z_{xx}-z_{yy}=4z_{uv}.
\]

\section{Graded Practice Problems}
\begin{enumerate}
\item Find
\[
\frac{\partial(u,v)}{\partial(x,y)}
\]
for
\[
u=3x+2y,\qquad v=x-y.
\]
\item If
\[
u=xy,\qquad v=x/y,
\]
find the Jacobian.
\item Verify the reciprocal Jacobian relation for
\[
u=x+y,\qquad v=x-y.
\]
\item Derive the polar formulas for \(f_x,f_y\).
\item Find the Jacobian of
\[
x=r\cos\theta,\qquad y=r\sin\theta.
\]
\end{enumerate}

% ============================================================
\chapter{Homogeneous Functions and Euler's Theorem}
% ============================================================

\section*{Learning Objectives}
Students should be able to:
\begin{itemize}
\item identify homogeneous functions and their degree;
\item state and prove Euler's theorem for two variables;
\item derive first- and second-order consequences;
\item apply Euler's theorem to composite functions efficiently.
\end{itemize}

\section*{Historical Motivation}

Homogeneous functions encode scaling. If all independent variables are
multiplied by the same factor, the output changes by a predictable
power of that factor. Such functions appear throughout geometry,
mechanics, thermodynamics and mathematical economics.

\section{Definition}

\begin{definition}
A function \(u=f(x,y)\) is homogeneous of degree \(n\) if
\[
f(tx,ty)=t^n f(x,y)
\]
for all suitable \(t>0\).
\end{definition}

\begin{example}
For
\[
f(x,y)=x^2+xy+y^2,
\]
\[
f(tx,ty)
=
t^2(x^2+xy+y^2).
\]
Hence \(f\) is homogeneous of degree \(2\).
\end{example}

\begin{example}
For
\[
f(x,y)=\frac{x+y}{x-y},
\]
\[
f(tx,ty)=f(x,y).
\]
Hence the degree is \(0\).
\end{example}

\section{Euler's Theorem}

\begin{theorem}[Euler's theorem for two variables]
Let
\[
u=f(x,y)
\]
be differentiable and homogeneous of degree \(n\). Then
\[
x\pd{u}{x}
+
y\pd{u}{y}
=
nu.
\]
\end{theorem}

\begin{proof}
Since \(u=f(x,y)\) is homogeneous of degree \(n\),
\[
f(tx,ty)=t^n f(x,y).
\]

Treat \(x,y\) as constants and differentiate both sides with respect
to \(t\).

By the chain rule,
\[
\begin{aligned}
\frac{d}{dt}f(tx,ty)
&=
f_x(tx,ty)\frac{d(tx)}{dt}
+
f_y(tx,ty)\frac{d(ty)}{dt}\\
&=
x f_x(tx,ty)
+
y f_y(tx,ty).
\end{aligned}
\]

On the right,
\[
\frac{d}{dt}
[t^n f(x,y)]
=
nt^{n-1}f(x,y).
\]

Therefore
\[
x f_x(tx,ty)
+
y f_y(tx,ty)
=
nt^{n-1}f(x,y).
\]

Put \(t=1\):
\[
x f_x(x,y)+y f_y(x,y)=nf(x,y).
\]
Since \(u=f(x,y)\),
\[
xu_x+yu_y=nu.
\]
\end{proof}

\section{First Differential Consequences}

Starting with
\[
xu_x+yu_y=nu,
\]
differentiate with respect to \(x\):
\[
u_x+xu_{xx}+yu_{yx}=nu_x.
\]
Assuming equality of the mixed derivatives under the appropriate
hypotheses,
\[
xu_{xx}+yu_{xy}=(n-1)u_x.
\]

Similarly, differentiation with respect to \(y\) gives
\[
xu_{xy}+yu_{yy}=(n-1)u_y.
\]

\section{Second-Order Euler Identity}

Multiply
\[
xu_{xx}+yu_{xy}=(n-1)u_x
\]
by \(x\):
\[
x^2u_{xx}+xyu_{xy}=(n-1)xu_x.
\]

Multiply
\[
xu_{xy}+yu_{yy}=(n-1)u_y
\]
by \(y\):
\[
xyu_{xy}+y^2u_{yy}=(n-1)yu_y.
\]

Add:
\[
x^2u_{xx}
+
2xyu_{xy}
+
y^2u_{yy}
=
(n-1)(xu_x+yu_y).
\]

By Euler's theorem,
\[
xu_x+yu_y=nu.
\]
Hence
\[
x^2u_{xx}
+
2xyu_{xy}
+
y^2u_{yy}
=
n(n-1)u.
\]

\section{Composite Form of Euler's Theorem}

Suppose \(v(x,y)\) is homogeneous of degree \(n\) and
\[
u=\phi(v).
\]
Then
\[
u_x=\phi'(v)v_x,
\qquad
u_y=\phi'(v)v_y.
\]
Therefore
\[
\begin{aligned}
xu_x+yu_y
&=
\phi'(v)(xv_x+yv_y)\\
&=
nv\phi'(v).
\end{aligned}
\]

This simple observation often saves a large amount of direct
differentiation.


\section{Recognizing Homogeneous Functions Efficiently}

A polynomial is homogeneous if every term has the same total degree.

For example,
\[
x^3+2x^2y-xy^2+5y^3
\]
is homogeneous of degree \(3\), because every monomial has total degree
\(3\).

A quotient
\[
\frac{P(x,y)}{Q(x,y)}
\]
of homogeneous functions of degrees \(p\) and \(q\) is homogeneous of
degree
\[
p-q,
\]
provided the quotient is defined.

Thus
\[
\frac{x^3+y^3}{x+y}
\]
is homogeneous of degree
\[
3-1=2.
\]

Similarly,
\[
(x^2+y^2)^{3/2}
\]
is homogeneous of degree \(3\), because
\[
[(tx)^2+(ty)^2]^{3/2}
=
t^3(x^2+y^2)^{3/2}
\]
for \(t>0\).

\section{Euler's Theorem as a Scaling Law}

The identity
\[
xu_x+yu_y=nu
\]
has a simple interpretation.

The vector
\[
(x,y)
\]
points radially away from the origin. Euler's theorem measures the
change of a homogeneous function when both variables are scaled
simultaneously.

If
\[
u(tx,ty)=t^nu(x,y),
\]
then increasing the scale \(t\) changes \(u\) at a rate determined by
its degree \(n\). Euler's differential identity is exactly the
infinitesimal form of that scaling law.

\section{More Composite Applications}

Suppose \(v(x,y)\) is homogeneous of degree \(n\) and
\[
u=\phi(v).
\]

Then
\[
u_x=\phi'(v)v_x,
\qquad
u_y=\phi'(v)v_y.
\]

Therefore
\[
xu_x+yu_y
=
\phi'(v)(xv_x+yv_y)
=
nv\phi'(v).
\]

This single formula generates many useful special cases.

If
\[
u=\log v,
\]
then
\[
\phi'(v)=\frac1v,
\]
so
\[
xu_x+yu_y=n.
\]

If
\[
u=e^v,
\]
then
\[
xu_x+yu_y=nve^v.
\]

If
\[
u=\sin^{-1}v,
\]
then
\[
xu_x+yu_y
=
\frac{nv}{\sqrt{1-v^2}}.
\]

Since \(v=\sin u\),
\[
xu_x+yu_y
=
n\tan u.
\]

If
\[
u=\tan^{-1}v,
\]
then
\[
xu_x+yu_y
=
\frac{nv}{1+v^2}.
\]

Since \(v=\tan u\),
\[
xu_x+yu_y
=
n\sin u\cos u
=
\frac n2\sin2u.
\]

\section{A Third-Order Pattern}

Euler's theorem can be differentiated repeatedly.

For a homogeneous function \(u\) of degree \(n\),
\[
xu_x+yu_y=nu.
\]

Differentiating once lowers the degree by one in the corresponding
derivative. Thus \(u_x\) and \(u_y\), when nonzero, are homogeneous of
degree \(n-1\).

Differentiating twice leads to the second-order identity
\[
x^2u_{xx}
+
2xyu_{xy}
+
y^2u_{yy}
=
n(n-1)u.
\]

The pattern suggests that repeated radial differentiation produces the
falling factorial
\[
n(n-1)(n-2)\cdots.
\]

This observation is useful when higher-order Euler identities are
encountered.

\section{Additional Detailed Solved Problems}

\subsection*{Problem 5. Verify Euler's theorem directly}

Let
\[
u=(x^2+y^2)^2.
\]

Since
\[
u(tx,ty)
=
[t^2(x^2+y^2)]^2
=
t^4u(x,y),
\]
the function is homogeneous of degree \(4\).

Now
\[
u_x
=
4x(x^2+y^2),
\]
and
\[
u_y
=
4y(x^2+y^2).
\]

Therefore
\[
\begin{aligned}
xu_x+yu_y
&=
4x^2(x^2+y^2)
+
4y^2(x^2+y^2)\\
&=
4(x^2+y^2)^2\\
&=
4u.
\end{aligned}
\]

This agrees with Euler's theorem.

\subsection*{Problem 6. Logarithm of a homogeneous function}

Let
\[
u
=
\log\left(\frac{x^2+y^2}{x+y}\right).
\]

The function
\[
v=
\frac{x^2+y^2}{x+y}
\]
is homogeneous of degree
\[
2-1=1.
\]

Since
\[
u=\log v,
\]
the composite Euler formula gives
\[
xu_x+yu_y=1.
\]

This avoids a lengthy direct differentiation.

\subsection*{Problem 7. Second-order identity}

If \(u\) is homogeneous of degree \(5\), simplify
\[
x^2u_{xx}+2xyu_{xy}+y^2u_{yy}.
\]

By the second-order Euler identity,
\[
x^2u_{xx}+2xyu_{xy}+y^2u_{yy}
=
n(n-1)u.
\]

Putting \(n=5\),
\[
x^2u_{xx}+2xyu_{xy}+y^2u_{yy}
=
20u.
\]



\section{Homogeneity: Several Equivalent Ways to Recognize It}

The defining condition is
\[
f(tx,ty)=t^nf(x,y).
\]

But in practice there are several useful recognition rules.

\subsection*{Polynomials}

A polynomial is homogeneous of degree \(n\) if every monomial has total
degree \(n\).

For example,
\[
x^4+3x^3y+2x^2y^2+y^4
\]
is homogeneous of degree \(4\).

\subsection*{Products}

If \(f\) is homogeneous of degree \(p\) and \(g\) is homogeneous of
degree \(q\), then
\[
fg
\]
is homogeneous of degree
\[
p+q.
\]

\subsection*{Quotients}

If \(g\ne0\),
\[
\frac{f}{g}
\]
is homogeneous of degree
\[
p-q.
\]

\subsection*{Powers}

If \(f\) is positive and homogeneous of degree \(p\), then
\[
f^\alpha
\]
is homogeneous of degree
\[
\alpha p.
\]

These rules often make the degree obvious without substituting
\((tx,ty)\) explicitly.

\section{Euler's Theorem: A Second Proof Viewpoint}

Define
\[
\phi(t)=f(tx,ty).
\]

Homogeneity gives
\[
\phi(t)=t^nf(x,y).
\]

Differentiate:
\[
\phi'(t)
=
x f_x(tx,ty)
+
y f_y(tx,ty).
\]

On the other hand,
\[
\phi'(t)
=
nt^{n-1}f(x,y).
\]

At \(t=1\),
\[
x f_x(x,y)+y f_y(x,y)=nf(x,y).
\]

This viewpoint is useful because Euler's theorem is simply the ordinary
one-variable derivative of the scaling function \(\phi(t)\).

\section{Converse of Euler's Theorem}

Under suitable regularity and on an appropriate star-shaped domain, the
Euler identity can also characterize homogeneity.

Suppose
\[
xf_x+yf_y=nf.
\]

Fix \(x,y\) and define
\[
\phi(t)=f(tx,ty).
\]

Then
\[
\phi'(t)
=
x f_x(tx,ty)+y f_y(tx,ty).
\]

Applying the Euler relation at \((tx,ty)\),
\[
(tx)f_x(tx,ty)
+
(ty)f_y(tx,ty)
=
n f(tx,ty).
\]

Therefore
\[
t\phi'(t)=n\phi(t),
\]
so
\[
\frac{\phi'(t)}{\phi(t)}
=
\frac nt
\]
where \(\phi\ne0\).

Integrating,
\[
\log|\phi(t)|
=
n\log t+C.
\]

Thus
\[
\phi(t)=Ct^n.
\]

At \(t=1\),
\[
C=\phi(1)=f(x,y).
\]

Hence
\[
f(tx,ty)=t^nf(x,y).
\]

This shows why Euler's differential equation and homogeneity are
closely connected.

\section{First Derivatives of a Homogeneous Function}

If \(u\) is homogeneous of degree \(n\), then \(u_x\) and \(u_y\) are
homogeneous of degree \(n-1\), provided the derivatives exist.

To see this, differentiate
\[
u(tx,ty)=t^nu(x,y)
\]
with respect to \(x\).

The left-hand side gives
\[
t\,u_x(tx,ty),
\]
while the right-hand side gives
\[
t^nu_x(x,y).
\]

Thus
\[
u_x(tx,ty)
=
t^{n-1}u_x(x,y).
\]

Therefore \(u_x\) is homogeneous of degree \(n-1\).

Similarly,
\[
u_y(tx,ty)
=
t^{n-1}u_y(x,y).
\]

This explains naturally why the differentiated Euler identities contain
the factor \(n-1\).

\section{Second Derivatives of a Homogeneous Function}

By the same reasoning,
\[
u_{xx},\quad u_{xy},\quad u_{yy}
\]
are homogeneous of degree
\[
n-2.
\]

Thus repeated differentiation lowers the degree by one each time.

\section{Deriving the Second-Order Euler Identity in Full}

Start from
\[
xu_x+yu_y=nu.
\]

Differentiate with respect to \(x\):
\[
u_x+xu_{xx}+yu_{yx}=nu_x.
\]

Hence
\[
xu_{xx}+yu_{xy}
=
(n-1)u_x.
\]

Differentiate the original Euler identity with respect to \(y\):
\[
xu_{xy}+u_y+yu_{yy}=nu_y.
\]

Hence
\[
xu_{xy}+yu_{yy}
=
(n-1)u_y.
\]

Multiply the first equation by \(x\):
\[
x^2u_{xx}+xyu_{xy}
=
(n-1)xu_x.
\]

Multiply the second by \(y\):
\[
xyu_{xy}+y^2u_{yy}
=
(n-1)yu_y.
\]

Add:
\[
x^2u_{xx}
+
2xyu_{xy}
+
y^2u_{yy}
=
(n-1)(xu_x+yu_y).
\]

Using
\[
xu_x+yu_y=nu,
\]
we obtain
\[
x^2u_{xx}
+
2xyu_{xy}
+
y^2u_{yy}
=
n(n-1)u.
\]

\section{Composite Homogeneous Functions: A General Template}

Let
\[
v=v(x,y)
\]
be homogeneous of degree \(n\), and let
\[
u=\phi(v).
\]

Then
\[
u_x=\phi'(v)v_x,
\qquad
u_y=\phi'(v)v_y.
\]

Therefore
\[
xu_x+yu_y
=
\phi'(v)(xv_x+yv_y).
\]

Euler's theorem for \(v\) gives
\[
xv_x+yv_y=nv.
\]

Hence
\[
xu_x+yu_y
=
nv\phi'(v).
\]

This one formula should be remembered as a method rather than as a
collection of separate examples.

\section{More Detailed Solved Problems}

\subsection*{Problem 8. Inverse trigonometric composite}

Let
\[
v=\frac{x^2+y^2}{x+y},
\]
and
\[
u=\sin^{-1}v.
\]

The numerator has degree \(2\) and the denominator degree \(1\), so
\(v\) is homogeneous of degree \(1\).

Therefore
\[
xv_x+yv_y=v.
\]

Now
\[
u_x
=
\frac{v_x}{\sqrt{1-v^2}},
\qquad
u_y
=
\frac{v_y}{\sqrt{1-v^2}}.
\]

Hence
\[
xu_x+yu_y
=
\frac{v}{\sqrt{1-v^2}}.
\]

Since
\[
v=\sin u,
\]
we obtain
\[
xu_x+yu_y
=
\frac{\sin u}{\cos u}
=
\tan u,
\]
where the relevant branch makes \(\cos u\) positive.

\subsection*{Problem 9. Exponential composite}

Let
\[
v=x^2+xy+y^2,
\qquad
u=e^v.
\]

Since \(v\) is homogeneous of degree \(2\),
\[
xv_x+yv_y=2v.
\]

Also,
\[
u_x=e^v v_x,
\qquad
u_y=e^v v_y.
\]

Therefore
\[
xu_x+yu_y
=
e^v(xv_x+yv_y)
=
2ve^v.
\]

Since
\[
u=e^v,
\]
this may also be written
\[
xu_x+yu_y=2vu.
\]

\subsection*{Problem 10. Second-order application}

Let
\[
u=x^4+2x^2y^2+y^4.
\]

This is homogeneous of degree \(4\).

Therefore, without calculating any second partial derivatives,
\[
x^2u_{xx}
+
2xyu_{xy}
+
y^2u_{yy}
=
4(3)u
=
12u.
\]

This is often much shorter than direct differentiation.


\section{Detailed Solved Problems}

\subsection*{Problem 1 (2025). Euler's theorem}
Let \(u=f(x,y)\) be homogeneous of degree \(n\). Show that
\[
xu_x+yu_y=nu.
\]

Since
\[
f(tx,ty)=t^n f(x,y),
\]
differentiate with respect to \(t\):
\[
x f_x(tx,ty)+y f_y(tx,ty)
=
nt^{n-1}f(x,y).
\]
Putting \(t=1\) gives
\[
xu_x+yu_y=nu.
\]

\subsection*{Problem 2 (2024). A trigonometric application}
Let
\[
u
=
\tan^{-1}
\left(
\frac{x^3+y^3}{x+y}
\right).
\]
Show that
\[
xu_x+yu_y=\sin2u.
\]

Set
\[
v=\frac{x^3+y^3}{x+y}.
\]
Since
\[
x^3+y^3=(x+y)(x^2-xy+y^2),
\]
we have
\[
v=x^2-xy+y^2,
\]
where the original expression is defined. Thus \(v\) is homogeneous
of degree \(2\).

Euler's theorem gives
\[
xv_x+yv_y=2v.
\]

Now
\[
u=\tan^{-1}v,
\]
so
\[
u_x=\frac{v_x}{1+v^2},
\qquad
u_y=\frac{v_y}{1+v^2}.
\]
Therefore
\[
\begin{aligned}
xu_x+yu_y
&=
\frac{xv_x+yv_y}{1+v^2}\\
&=
\frac{2v}{1+v^2}.
\end{aligned}
\]

Since
\[
v=\tan u,
\]
\[
\begin{aligned}
xu_x+yu_y
&=
\frac{2\tan u}{1+\tan^2u}\\
&=
2\tan u\cos^2u\\
&=
2\sin u\cos u\\
&=
\sin2u.
\end{aligned}
\]

\subsection*{Problem 3 (2023). A differentiated Euler identity}
Suppose \(u(x,y)\) is homogeneous of degree \(n\). Simplify
\[
xu_{xy}+yu_{yy}.
\]

Euler's theorem gives
\[
xu_x+yu_y=nu.
\]
Differentiate with respect to \(y\):
\[
xu_{xy}+u_y+yu_{yy}=nu_y.
\]
Hence
\[
xu_{xy}+yu_{yy}=(n-1)u_y.
\]

\subsection*{Problem 4. Logarithm of a homogeneous function}
Suppose \(v(x,y)\) is homogeneous of degree \(n\), and
\[
u=\log v.
\]
Then
\[
u_x=\frac{v_x}{v},
\qquad
u_y=\frac{v_y}{v}.
\]
Thus
\[
xu_x+yu_y
=
\frac{xv_x+yv_y}{v}
=
\frac{nv}{v}
=
n.
\]

\section{Graded Practice Problems}

\subsection*{Level I}
\begin{enumerate}
\item Determine the degree of
\[
x^3+x^2y+xy^2+y^3.
\]
\item Determine the degree of
\[
\frac{x^2+y^2}{x+y}.
\]
\item Verify Euler's theorem for
\[
u=x^2+xy+y^2.
\]
\end{enumerate}

\subsection*{Level II}
\begin{enumerate}
\item Verify Euler's theorem for
\[
u=(x^2+y^2)^{3/2}.
\]
\item If \(u\) is homogeneous of degree \(n\), prove
\[
xu_{xx}+yu_{xy}=(n-1)u_x.
\]
\item Derive
\[
x^2u_{xx}+2xyu_{xy}+y^2u_{yy}=n(n-1)u.
\]
\end{enumerate}

\subsection*{Level III}
\begin{enumerate}
\item If \(v\) is homogeneous of degree \(n\) and \(u=\log v\),
show that \(xu_x+yu_y=n\).
\item If \(v\) is homogeneous of degree \(n\) and \(u=e^v\), show
that
\[
xu_x+yu_y=nve^v.
\]
\item If \(v\) is homogeneous of degree \(n\) and
\(u=\sin^{-1}v\), express \(xu_x+yu_y\) in terms of \(u\).
\end{enumerate}

% ============================================================
\chapter{Integrated Review and Problem-Solving Guide}
% ============================================================

\section{How to Approach a Two-Variable Limit}

When asked to examine
\[
\lim_{(x,y)\to(a,b)}f(x,y),
\]
use the following sequence.

\begin{enumerate}
\item Try direct substitution.
\item If an indeterminate form occurs, simplify algebraically.
\item If the point is the origin and \(x^2+y^2\) is present, consider
polar coordinates.
\item To prove nonexistence, test simple paths:
\[
x=0,\qquad y=0,\qquad y=mx.
\]
\item If these paths agree, do not conclude existence. Consider curved
paths such as \(y=x^2\) when suggested by the algebra.
\item To prove existence, use an estimate, the Squeeze Theorem, polar
coordinates, or the \(\eps\)-\(\delta\) definition.
\end{enumerate}

\section{How to Test Differentiability}

At \((a,b)\):

\begin{enumerate}
\item Check continuity first. A discontinuous function cannot be
differentiable.
\item Find \(f_x(a,b)\) and \(f_y(a,b)\).
\item If \(f_x,f_y\) exist in a neighbourhood and are continuous at
the point, differentiability follows.
\item Otherwise test
\[
\frac{
f(a+h,b+k)-f(a,b)-f_x(a,b)h-f_y(a,b)k
}{
\sqrt{h^2+k^2}
}
\to0.
\]
\end{enumerate}

\section{Common Misconceptions}

\begin{enumerate}
\item Equal values along the coordinate axes do not prove a joint
limit exists.
\item Agreement along all straight lines still need not prove a limit
exists.
\item Equal repeated limits do not guarantee the joint limit.
\item Existence of \(f_x,f_y\) does not imply continuity.
\item Existence of \(f_x,f_y\) does not imply differentiability.
\item Mixed partial derivatives are not automatically equal without
suitable hypotheses.
\item Euler's theorem applies only after homogeneity and its degree
have been identified correctly.
\item In a chain-rule problem, every path of dependence contributes a
term.
\end{enumerate}

\section{Quick Reference}

\subsection*{Joint limit}
\[
\lim_{(x,y)\to(a,b)}f(x,y)=L
\]
means: for every \(\eps>0\), there exists \(\delta>0\) such that
\[
0<\sqrt{(x-a)^2+(y-b)^2}<\delta
\]
implies
\[
\abs{f(x,y)-L}<\eps.
\]

\subsection*{Continuity}
\[
f\text{ continuous at }(a,b)
\iff
\lim_{(x,y)\to(a,b)}f(x,y)=f(a,b).
\]

\subsection*{Partial derivatives}
\[
f_x(a,b)
=
\lim_{h\to0}
\frac{f(a+h,b)-f(a,b)}{h},
\]
\[
f_y(a,b)
=
\lim_{k\to0}
\frac{f(a,b+k)-f(a,b)}{k}.
\]

\subsection*{Differentiability}
\[
f(a+h,b+k)-f(a,b)
=
f_x(a,b)h+f_y(a,b)k
+
o(\sqrt{h^2+k^2}).
\]

\subsection*{Tangent plane}
\[
z=f(a,b)+f_x(a,b)(x-a)+f_y(a,b)(y-b).
\]

\subsection*{Chain rule}
If
\[
z=f(u,v),\quad u=u(x,y),\quad v=v(x,y),
\]
then
\[
z_x=f_u u_x+f_v v_x,
\qquad
z_y=f_u u_y+f_v v_y.
\]

\subsection*{Jacobian}
\[
\frac{\partial(u,v)}{\partial(x,y)}
=
\begin{vmatrix}
u_x&u_y\\
v_x&v_y
\end{vmatrix}.
\]

\subsection*{Euler's theorem}
If \(u\) is homogeneous of degree \(n\),
\[
xu_x+yu_y=nu.
\]
Also,
\[
xu_{xx}+yu_{xy}=(n-1)u_x,
\]
\[
xu_{xy}+yu_{yy}=(n-1)u_y,
\]
and
\[
x^2u_{xx}+2xyu_{xy}+y^2u_{yy}=n(n-1)u.
\]


\section{Conceptual Checklist}

Before finishing the unit, a student should be able to answer the
following questions without relying only on memorized formulas.

\begin{enumerate}
\item Why is a two-variable limit harder than a one-variable limit?
\item Why can two different paths disprove a limit?
\item Why can checking many paths still fail to prove a limit?
\item What is the difference between a joint limit and a repeated
limit?
\item Why do partial derivatives examine only coordinate directions?
\item Why does existence of partial derivatives not guarantee
differentiability?
\item What does the error term in the definition of differentiability
mean?
\item How is the tangent plane obtained from the linear part of the
increment?
\item Under what hypotheses may we interchange the order of mixed
partial differentiation?
\item How does the chain rule follow dependency paths?
\item What does a Jacobian measure algebraically and geometrically?
\item Why does Euler's theorem express the scaling property of a
homogeneous function?
\end{enumerate}

\section{Further Comprehensive Problems}

\begin{enumerate}
\item Find and sketch the domain of
\[
f(x,y)=\log(9-x^2-y^2)+\sqrt{y-x^2}.
\]

\item Determine whether
\[
\lim_{(x,y)\to(0,0)}
\frac{x^2y^2}{x^4+y^4}
\]
exists.

\item Show that
\[
\lim_{(x,y)\to(0,0)}
\frac{x^4y^2}{x^4+y^2}=0.
\]

\item Find both repeated limits of
\[
f(x,y)=\frac{x^2-y^2}{x^2+y^2}
\]
at the origin and compare them with the joint limit.

\item Prove directly from the definition that
\[
f(x,y)=x^2+xy+y^2
\]
is differentiable at an arbitrary point \((a,b)\).

\item Find the tangent plane to
\[
z=x^2e^y
\]
at \((1,0)\).

\item For
\[
z=u^2+uv+v^2,\quad
u=x+y,\quad
v=x-y,
\]
find \(z_x,z_y\) by the chain rule and verify the result by first
expressing \(z\) directly in terms of \(x,y\).

\item For
\[
u=x^2-y^2,\qquad v=2xy,
\]
calculate the Jacobian and determine where the transformation can fail
to be locally invertible.

\item If \(u(x,y)\) is homogeneous of degree \(n\), derive
\[
xu_{xx}+yu_{xy}=(n-1)u_x
\]
and
\[
xu_{xy}+yu_{yy}=(n-1)u_y.
\]

\item Let
\[
v=x^2+xy+y^2,\qquad
u=\tan^{-1}v.
\]
Use Euler's theorem to simplify
\[
xu_x+yu_y.
\]
\end{enumerate}


\section{Comprehensive Review Problems}

\begin{enumerate}
\item Examine
\[
\lim_{(x,y)\to(0,0)}
\frac{x^2y}{x^4+y^2}.
\]

\item Determine both repeated limits of
\[
\frac{x-y}{x+y}.
\]

\item Find all first- and second-order partial derivatives of
\[
u=e^{xy}\sin(x+y).
\]

\item Determine whether
\[
f(x,y)=
\begin{cases}
\dfrac{x^2y}{x^2+y^2},&(x,y)\ne(0,0),\\
0,&(0,0)
\end{cases}
\]
is continuous at the origin.

\item Explain why existence of both first partial derivatives at a
point does not guarantee differentiability.

\item Let
\[
z=u^2v,\quad u=x+y,\quad v=x-y.
\]
Find \(z_x,z_y\).

\item For
\[
u=x^2-y^2,\qquad v=2xy,
\]
calculate
\[
\frac{\partial(u,v)}{\partial(x,y)}.
\]

\item If \(u\) is homogeneous of degree \(5\), simplify
\[
xu_x+yu_y.
\]

\item If \(u\) is homogeneous of degree \(n\), simplify
\[
xu_{xy}+yu_{yy}.
\]

\item If
\[
v=x^2+xy+y^2,
\qquad
u=\log v,
\]
find \(xu_x+yu_y\) without separately expanding both partial
derivatives.
\end{enumerate}

\backmatter

\chapter*{Final Perspective}
\addcontentsline{toc}{chapter}{Final Perspective}

The essential change from one-variable to multivariable calculus is
geometrical. A point on a line has only two basic directions of
approach, while a point in the plane can be approached along infinitely
many curves. A one-variable derivative gives one local slope; a
two-variable function has several directional rates of change.

Partial derivatives isolate coordinate directions, while
differentiability brings all small changes together through a single
linear approximation. The chain rule explains how changes propagate
through intermediate variables. Jacobians organize changes of
coordinates, and Euler's theorem converts a scaling law into a
differential identity.

These ideas are foundational for vector calculus, partial differential
equations, differential geometry, mathematical physics and applied
mathematics.

\end{document}
